\chapter{数学工具与物理约定}\label{chap:math-foundations}

普通量子力学已经提供三件会反复使用的工具：动量表象、Hilbert 空间和 Green 函数。进入相对论场论以后，这些工具仍然有效，但必须同时服从四维度规、连续谱归一化和统一的 Fourier 约定。Fourier 变换的正负号、Minkowski 度规、Green 函数的边界条件与连续态测度会进入同一条计算；只要其中一项前后不一致，$2\pi$、$\hbar$ 或整体符号就会随之出错。

这些问题自然给出一条由熟悉对象通向场论的路线：单位、坐标和 Fourier 约定先固定“怎样表示”；连续谱、张量积和投影再固定“量子态怎样组织”；有限维 Gaussian 积分提供函数积分的代数原型；参数积分与长时间极限则把这些工具接到实际散射计算。

\section{单位、Fourier 变换与分布}

Fourier 变换把微分方程化成代数方程，也把时空局域性转化成能量—动量守恒；因此传播子、谱密度和散射振幅都会反复依赖同一套号规。先固定单位、度规和 Fourier 变换，Dirac $\delta$、极点边界值与连续动量测度便随之唯一确定。

全文固定使用 SI 制并显式保留 $c,\hbar,\varepsilon_0,\mu_0,k_B$。Minkowski 度规取
\begin{equation}
\eta_{\mu\nu}=\operatorname{diag}(1,-1,-1,-1),
\qquad x^\mu=(ct,\mathbf r),
\qquad p^\mu=(E/c,\mathbf p).
\label{eq:global-metric-convention}
\end{equation}
希腊指标 $\mu,\nu=0,1,2,3$，拉丁空间指标 $i,j=1,2,3$；重复 Lorentz 指标求和。Fourier 四波矢记作 $\kappa^\mu=(\omega/c,\boldsymbol\kappa)$，真实四动量则是 $p^\mu$；特别地，对光子 $k^\mu=\hbar\kappa^\mu$。这个区分避免把“波矢”和“动量”混为同一个量纲。

时间 Fourier 对固定为
\begin{equation}
\widetilde f(\omega)=\int_{-\infty}^{\infty}\dd t\,f(t)\ee^{+\ii\omega t},
\qquad
f(t)=\int_{-\infty}^{\infty}\frac{\dd\omega}{2\pi}\,\widetilde f(\omega)\ee^{-\ii\omega t}.
\label{eq:global-time-fourier}
\end{equation}
空间波矢与物理动量两种表示分别使用 $\dd^3 \kappa/(2\pi)^3$ 和 $\dd^3 p/(2\pi\hbar)^3$。Fourier 反演也固定 $\delta$ 分布、$\Theta$ 函数、$\ii0^+$ 处方和积分轮廓边界值之间的关系。

% -----------------------------------------------------------------------------
% Fourier, distribution and contour-analysis foundations before propagators
% -----------------------------------------------------------------------------

传播子的定义反复依赖 Fourier 变换、分布与复分析中的边界值，因此这些数学约定必须先由定义固定。这些工具按依赖关系自然出现：测试函数先给出 Dirac $\delta$ 的定义，主值与 $\ii0^+$ 边界值随后规定奇点怎样取极限，复能量积分的闭合方向最后把这些边界值转化为不同传播条件。后续“推迟”“超前”或“Feynman”传播子只是在这些数学对象上加入不同物理边界条件。

Dirac $\delta$ 不是普通函数，而是作用在测试函数上的线性泛函。对连续测试函数 $f$，其平移形式的定义是
\begin{equation}
 \boxed{
 \int\dd y\,\delta(x-y)f(y)=f(x).}
 \label{eq:delta-def-action-r13}
\end{equation}
与全书 Fourier 号规相容的常用表示为
\begin{equation}
 \boxed{
 \delta(t-t')
 =\int\frac{\dd\omega}{2\pi}\ee^{-\ii\omega(t-t')},
 \qquad
 \delta^{(3)}(\mathbf r-\mathbf r')
 =\int\frac{\dd^3 p}{(2\pi\hbar)^3}
 \ee^{+\ii\mathbf p\cdot(\mathbf r-\mathbf r')/\hbar}.}
 \label{eq:delta-fourier-representation-dp68}
\end{equation}
质量壳 $\delta$ 函数和长时间 $\delta$ 极限都应通过测试函数作用或变量替换理解，而不能把 $\delta$ 当成在一点取无穷大的普通函数。

Heaviside 分布的导数由分布导数定义固定：
\begin{equation}
 \boxed{\frac{\dd\Theta(x)}{\dd x}=\delta(x).}
 \label{eq:heaviside-delta-dp68}
\end{equation}
因此时间排序对象求导时会产生等时接触项。若 $F(t)$ 在 $t=t'$ 连续可微，则
\begin{equation}
 \boxed{
 \partial_t[\Theta(t-t')F(t)]
 =\delta(t-t')F(t')
 +\Theta(t-t')\partial_tF(t).}
 \label{eq:theta-product-derivative-dp68}
\end{equation}

复平面边界值从有限 $\epsilon>0$ 开始。对实变量 $x$，
\begin{equation}
 \boxed{
 \frac1{x\pm\ii\epsilon}
 =\frac{x}{x^2+\epsilon^2}
 \mp\ii\frac{\epsilon}{x^2+\epsilon^2}.}
 \label{eq:plemelj-split-r13}
\end{equation}
这里 $\epsilon$ 与 $x$ 具有相同量纲。主值分布定义为
\begin{equation}
 \boxed{
 \left\langle\operatorname{PV}\frac1x,f\right\rangle
 \equiv\lim_{\epsilon\to0^+}
 \left(\int_{-\infty}^{-\epsilon}+\int_{\epsilon}^{\infty}\right)
 \frac{f(x)}x\,\dd x.}
 \label{eq:principal-value-definition-dp68}
\end{equation}
在分布意义下取 $\epsilon\to0^+$，正文\eqnrefs{eq:plemelj-split-r13} 变成
\begin{equation}
 \boxed{
 \frac1{x\pm\ii0^+}
 =\operatorname{PV}\frac1x\mp\ii\pi\delta(x).}
 \label{eq:sokhotski-plemelj-dp68}
\end{equation}
这条边界值恒等式直接控制传播子的虚部与在壳 $\delta$ 项。

轮廓闭合方向由 Fourier 指数在大半圆上的衰减决定。令 $z=\operatorname{Re}z+\ii\operatorname{Im}z$，取实时间差 $\tau$，则
\begin{equation}
 \boxed{
 \left|\ee^{-\ii z\tau/\hbar}\right|
 =\exp\!\left(\frac{\tau\,\operatorname{Im}z}{\hbar}\right).}
 \label{eq:jordan-exponential-r13}
\end{equation}
因此 $\tau>0$ 时向下闭合、$\tau<0$ 时向上闭合。若
$I(\tau)=\int_{-\infty}^{\infty}\dd z\,F(z)\ee^{-\ii z\tau/\hbar}/(2\pi)$
且半圆积分消失，则取向进一步给出
\begin{equation}
 \boxed{
 \begin{aligned}
 \tau>0:&\quad
 I(\tau)=-\ii\sum_{z_j\in\mathbb C_-}
 \operatorname{Res}_{z_j}[F(z)\ee^{-\ii z\tau/\hbar}],\\
 \tau<0:&\quad
 I(\tau)=+\ii\sum_{z_j\in\mathbb C_+}
 \operatorname{Res}_{z_j}[F(z)\ee^{-\ii z\tau/\hbar}].
 \end{aligned}}
 \label{eq:contour-orientation-residue-dp68}
\end{equation}
下半圆是顺时针，因此第一行多一个负号；这与“选哪些极点”是两个不同的问题。

若 $z_0$ 是 $m$ 阶极点，可局部写成
\begin{equation}
 \boxed{F(z)=\frac{g(z)}{(z-z_0)^m},\qquad g(z)\ \text{在 }z_0\text{ 解析}.}
 \label{eq:higher-pole-local-form-dp68}
\end{equation}
其留数为
\begin{equation}
 \boxed{
 \operatorname{Res}_{z=z_0}F(z)
 =\frac1{(m-1)!}
 \left.\frac{\dd^{m-1}}{\dd z^{m-1}}
 [(z-z_0)^mF(z)]\right|_{z=z_0}.}
 \label{eq:higher-pole-residue-dp68}
\end{equation}

% DeepPass21: detailed Fourier/distribution derivations.

\begin{equation}
\boxed{
\widetilde{\partial_t f}(\omega)=-\ii\omega\widetilde f(\omega),
\qquad
\widetilde{f*g}(\omega)=\widetilde f(\omega)\widetilde g(\omega).}
\label{eq:fourier-derivative-convolution-dp21}
\end{equation}

\begin{equation}
\boxed{
\int_{-\infty}^{\infty}\dd t\,f^*(t)g(t)
=\int_{-\infty}^{\infty}\frac{\dd\omega}{2\pi}\,
\widetilde f^{\,*}(\omega)\widetilde g(\omega).}
\label{eq:parseval-time-dp21}
\end{equation}




\begin{derivation}[从\eqnrefs{eq:plemelj-split-r13}到\eqnrefs{eq:sokhotski-plemelj-dp68}]

正文\eqnrefs{eq:plemelj-split-r13} 已把复分母分成实部和虚部两个普通函数。$\epsilon\to0^+$ 时，实部按正文\eqnrefs{eq:principal-value-definition-dp68} 趋于 Cauchy 主值，虚部则形成单位面积、宽度趋零的 Lorentz 峰；分别取这两个分布极限即可得到正文\eqnrefs{eq:sokhotski-plemelj-dp68}。

先令
\begin{equation*}
 \delta_\epsilon(x)
 =\frac1\pi\frac{\epsilon}{x^2+\epsilon^2}.
\end{equation*}
变量代换 $x=\epsilon u$ 给
\begin{equation*}
 \int\dd x\,\delta_\epsilon(x)=1,
\end{equation*}
且对任意光滑测试函数 $f$，
\begin{equation*}
 \int\dd x\,\delta_\epsilon(x)f(x)
 =\frac1\pi\int\frac{\dd u}{1+u^2}f(\epsilon u)
 \longrightarrow f(0).
\end{equation*}
所以 $\delta_\epsilon\to\delta$。另一方面，把
$x/(x^2+\epsilon^2)$ 的正负半轴配对后，极限正是正文\eqnrefs{eq:principal-value-definition-dp68}。代回正文\eqnrefs{eq:plemelj-split-r13} 即得正文\eqnrefs{eq:sokhotski-plemelj-dp68}。
\end{derivation}

\begin{derivation}[从\eqnrefs{eq:jordan-exponential-r13}到\eqnrefs{eq:contour-orientation-residue-dp68}]

正文\eqnrefs{eq:jordan-exponential-r13} 已告诉我们指数在复平面的哪一侧衰减。还需要再检查大半圆积分确实消失以及闭合曲线的取向：$\tau>0$ 时下半圆顺时针，$\tau<0$ 时上半圆逆时针。两点合起来便得到正文\eqnrefs{eq:contour-orientation-residue-dp68} 的正负号。

设大 $|z|$ 时 $|F(z)|\le C/|z|^{1+\delta}$，$\delta>0$。对 $\tau>0$ 取下半圆 $z=Re^{-\ii\theta}$，则
\begin{equation*}
 |\ee^{-\ii z\tau/\hbar}|=\ee^{-R\tau\sin\theta/\hbar}\le1,
\end{equation*}
从而
\begin{equation*}
 \left|\int_{C_R^-}\dd z\,F(z)\ee^{-\ii z\tau/\hbar}\right|
 \le\frac{\pi C}{R^\delta}\to0.
\end{equation*}
下半圆从 $+R$ 绕回 $-R$，闭合曲线顺时针，因此留数定理给实轴积分
$-2\pi\ii$ 乘下半平面留数；除以正文 Fourier convention 中的 $2\pi$ 即得\eqnrefs{eq:contour-orientation-residue-dp68} 第一行。$\tau<0$ 时同理向上闭合，曲线逆时针，得到第二行。
\end{derivation}



\section{Hilbert 空间与投影}

Fourier 约定解决了“同一个对象怎样换表示”，却没有解决“哪些量子态属于我们要保留的系统”。因此下一步从表示空间转到状态空间，并用投影算符把保留自由度与被消去自由度严格区分。

\subsection{Hilbert 空间与连续态}
离散能级问题里，读者习惯把态写成一列可平方求和的系数；自由粒子却以连续动量标记本征态，归一化中必然出现 Dirac $\delta$ 函数。进入散射理论以前，必须先把“普通 Hilbert 空间向量”和“连续谱广义本征态”区分清楚，否则完备关系、态密度和相空间测度会混入彼此不兼容的归一化因子。

量子态属于 Hilbert 空间 $\mathcal H$；可观测量由定义域适当的自伴算符表示。若复合系统由 $A$ 与 $B$ 构成，则状态空间是张量积 $\mathcal H_A\otimes\mathcal H_B$，局域算符写成 $A\otimes I_B$ 与 $I_A\otimes B$。对两个算符定义
\begin{equation*}
[A,B]\equiv AB-BA,
\qquad
\{A,B\}\equiv AB+BA.
\end{equation*}
幺正算符满足 $U^\dagger U=I$；由自伴生成元 $G$ 产生的一参数连续变换写为 $U(s)=\ee^{-\ii sG/\hbar}$。Poincar\'e 群、时间演化和量子散射均采用这一约定。

连续本征态不是普通 Hilbert 空间向量，而是装备 Hilbert 空间意义下的广义态。连续谱常用两种等价但不能混写的归一化。这里分别定义它们。Fourier 归一化平面波基记作 $|\mathbf p,s\rangle_F$：
\begin{align}
{}_F\!\langle\mathbf p,s|\mathbf p',s'\rangle_F
&=(2\pi\hbar)^3\delta_{ss'}\delta^{(3)}(\mathbf p-\mathbf p'),\\
I&=\sum_s\int\frac{\dd^3 p}{(2\pi\hbar)^3}
|\mathbf p,s\rangle_F{}_F\!\langle\mathbf p,s|.
\label{eq:global-continuum-normalization-fourier}
\end{align}
自由电子光学的连续积分核更适合使用纯 Dirac-$\delta$ 归一化，定义
\begin{equation}
|\mathbf p,s\rangle_\delta
\equiv(2\pi\hbar)^{-3/2}|\mathbf p,s\rangle_F.
\label{eq:delta-normalized-state-definition}
\end{equation}
逐项代入上式便得到
\begin{align}
{}_{\delta}\!\langle\mathbf p,s|\mathbf p',s'\rangle_{\delta}
&=\delta_{ss'}\delta^{(3)}(\mathbf p-\mathbf p'),\\
I&=\sum_s\int \dd^3 p\,
|\mathbf p,s\rangle_\delta{}_{\delta}\!\langle\mathbf p,s|.
\label{eq:global-continuum-normalization-delta}
\end{align}
若省略下标而直接写 $|\mathbf p,s\rangle$，在自由电子连续核章节默认指第二套 $\delta$ 归一化；相对论散射中使用的协变量子态归一化还会另外引入 $2E_{\mathbf p}$ 因子，并在相应章节从 Lorentz 不变测度推出。三套归一化的换算必须显式进行，不能通过“吸收常数”混用。

% DeepPass21: continuum normalization from a finite box.

把无限空间暂时放进边长为 $L$ 的周期立方盒，可以把连续动量谱变成离散格点。周期边界条件要求
\begin{equation}
 \boxed{
 p_i=\frac{2\pi\hbar}{L}n_i,
 \qquad n_i\in\mathbb Z,
 \qquad
 \Delta p_i=\frac{2\pi\hbar}{L}.}
 \label{eq:periodic-box-momenta-dp68}
\end{equation}
因此每个允许动量点在三维动量空间中占据体积 $(2\pi\hbar/L)^3$。当 $L\to\infty$ 时，离散 Riemann 和与位置完备核分别趋于
\begin{equation}
\boxed{
\sum_{\mathbf n}\longrightarrow
\frac{L^3}{(2\pi\hbar)^3}\int\dd^3 p,
\qquad
\frac1{L^3}\sum_{\mathbf n}\ee^{\ii\mathbf p_{\mathbf n}\cdot(\mathbf r-\mathbf r')/\hbar}
\longrightarrow\delta^{(3)}(\mathbf r-\mathbf r').}
\label{eq:box-to-continuum-dp21}
\end{equation}
这里第二个极限按分布意义理解；它不是逐点函数极限。

\begin{derivation}[从\eqnrefs{eq:periodic-box-momenta-dp68}到\eqnrefs{eq:box-to-continuum-dp21}]
对随动量变化足够平滑的测试函数 $F(\mathbf p)$，

\begin{align*}
 \sum_{\mathbf n}F(\mathbf p_{\mathbf n})
 &=\sum_{\mathbf n}
 \left(\frac{2\pi\hbar}{L}\right)^3
 \frac{F(\mathbf p_{\mathbf n})}{(2\pi\hbar/L)^3}\\
 &\longrightarrow
 \frac{L^3}{(2\pi\hbar)^3}
 \int\dd^3 p\,F(\mathbf p).
\end{align*}
这给出第一条极限。

盒归一化平面波为
\begin{equation*}
 \langle\mathbf r|\mathbf n\rangle
 =L^{-3/2}\ee^{\ii\mathbf p_{\mathbf n}\cdot\mathbf r/\hbar},
\end{equation*}
离散完备性
$\sum_{\mathbf n}|\mathbf n\rangle\langle\mathbf n|=I$
在位置表象中变成
\begin{equation*}
 \frac1{L^3}\sum_{\mathbf n}
 \ee^{\ii\mathbf p_{\mathbf n}\cdot(\mathbf r-\mathbf r')/\hbar}
 =\langle\mathbf r|I|\mathbf r'\rangle.
\end{equation*}
将刚得到的求和--积分极限用于左端，
\begin{equation*}
 \frac1{L^3}\sum_{\mathbf n}
 \ee^{\ii\mathbf p_{\mathbf n}\cdot\Delta\mathbf r/\hbar}
 \longrightarrow
 \int\frac{\dd^3 p}{(2\pi\hbar)^3}
 \ee^{\ii\mathbf p\cdot\Delta\mathbf r/\hbar}
 =\delta^{(3)}(\Delta\mathbf r),
\end{equation*}
即正文\eqnrefs{eq:box-to-continuum-dp21} 的第二条关系。
\end{derivation}




\subsection{投影与有效哈密顿量}
“保留哪些自由度”可以由投影算符精确定义。令 $P=P^2=P^\dagger$ 投影到目标子空间，$Q=I-P$；任意哈密顿量都可按 $P\oplus Q$ 分块。有效哈密顿量应由完整预解式的 Schur 补得到，而不是预先删去 $Q$ 空间，因为被消去自由度仍通过能量依赖的自能反馈到目标子空间。

% General projection/resolvent mathematics

许多有效理论都从同一个线性代数问题开始：完整 Hilbert 空间可分成希望显式保留的目标子空间 $P\mathcal H$ 与被消去的互补子空间 $Q\mathcal H$。定义正交投影
\begin{equation*}
 P^2=P=P^\dagger,
 \qquad
 Q\equiv I-P,
 \qquad
 PQ=QP=0.
\end{equation*}
完整 Hamilton 算符记为 $H$，单位为 J；复能量参数 $z$ 也以 J 为单位，并取在 $QHQ$ 的预解集中，使 $(z-QHQ)^{-1}$ 存在。

把 $z-H$ 按 $P\oplus Q$ 分块，先写成
\begin{equation}
 \boxed{
 z-H=
 \begin{pmatrix}
 z-PHP & -PHQ\\
 -QHP & z-QHQ
 \end{pmatrix}.}
 \label{eq:feshbach-block-matrix-dp68}
\end{equation}
对右下块作 Schur 补，目标子空间中的精确预解式为
\begin{equation}
 \boxed{
 P(z-H)^{-1}P=
 \left[z-PHP-PHQ(z-QHQ)^{-1}QHP\right]^{-1}.}
 \label{eq:feshbach-schur-proof-r10}
\end{equation}
这条式子没有假设 $P$--$Q$ 耦合很弱；它只是分块逆矩阵恒等式。由右端自然定义能量依赖的投影自能和精确有效 Hamilton 算符
\begin{equation}
 \boxed{
 H_{\rm eff}(z)
 \equiv PHP+\Sigma_P(z),
 \qquad
 \Sigma_P(z)
 \equiv PHQ(z-QHQ)^{-1}QHP.}
 \label{eq:feshbach_schur_exact_positive_projection}
\end{equation}
$H_{\rm eff}$ 与 $\Sigma_P$ 的单位均为 J。$\Sigma_P(z)$ 记录从 $P$ 子空间跃迁到 $Q$、在 $Q$ 中传播、再返回 $P$ 的全部线性虚过程；其能量依赖来自被消去子空间的预解式，不能在没有额外尺度条件时删去。

若目标能窗与 $QHQ$ 的谱之间存在有限能隙，能量依赖的精确预解式可在该谱隔离控制下展开为静态分块对角化。把
\begin{equation*}
 H=H_{\rm d}+H_{\rm od},
 \qquad
 H_{\rm d}=PHP+QHQ,
 \qquad
 H_{\rm od}=PHQ+QHP
\end{equation*}
分成分块对角与分块非对角部分，并取反厄米、纯非对角生成元 $S$，使 $\widetilde H=e^SHe^{-S}$。条件 $S^\dagger=-S$ 保证 $e^S$ 幺正，因此这一分块对角化不会改变内积或 Hamilton 算符的厄米性。为消去一阶非对角块，$S$ 必须满足 Sylvester 方程
\begin{equation}
 \boxed{H_{\rm od}+[S,H_{\rm d}]=0.}
 \label{eq:sw-sylvester-r15}
\end{equation}
Sylvester 方程固定 $S$ 的一阶非对角部分；同一个 $S$ 再通过 BCH 的二阶对易子产生目标子空间内的能级位移与诱导耦合。把 Baker--Campbell--Hausdorff 级数保留到 $S^2$，并使用\eqnrefs{eq:sw-sylvester-r15} 消去一阶非对角项后，得到
\begin{equation}
 \boxed{
 \widetilde H
 =H_{\rm d}+\frac12[S,H_{\rm od}]
 +\mathcal O(S^3H).}
 \label{eq:sw-bch-second-r15}
\end{equation}
$1/2$ 来自一阶对易子 $[S,H_{\rm od}]$ 与二阶嵌套对易子 $[S,[S,H_{\rm d}]]/2$ 在使用 Sylvester 方程后的系数组合。

进一步令 $|p\rangle$、$|p'\rangle$ 是 $P$ 子空间中的 $H_{\rm d}$ 本征态，$|q\rangle$ 是 $Q$ 子空间本征态，分别满足能量 $E_p,E_{p'},E_q$。二阶静态 Schrieffer--Wolff 有效 Hamilton 算符的矩阵元为
\begin{equation}
 \boxed{
 \begin{aligned}
 \langle p|H_{\rm eff}^{\rm SW}|p'\rangle
 ={}&E_p\delta_{pp'}\\
 &+\frac12\sum_{q\in Q}H_{pq}H_{qp'}
 \left[
 \frac1{E_p-E_q}
 +\frac1{E_{p'}-E_q}
 \right]
 +\mathcal O(\epsilon_{\rm SW}^3\Delta_Q).
 \end{aligned}}
 \label{eq:sw_arbitrary_target_subspace_matrix_elements}
\end{equation}
其中 $H_{pq}=\langle p|H|q\rangle$，$\Delta_Q$ 是目标能窗到 $QHQ$ 谱的最小能隙，
$\epsilon_{\rm SW}\equiv\|PHQ\|/\Delta_Q$ 是无量纲混合参数。只有 $\epsilon_{\rm SW}\ll1$ 时，上式的二阶截断才受控；否则应回到\eqnrefs{eq:feshbach_schur_exact_positive_projection} 的精确能量依赖形式。

% DeepPass21: detailed projection/resolvent identities.

对自伴算符 $A=QHQ$，谱定理给出
\begin{equation}
 \boxed{
 A=\int_{\sigma(A)}\lambda\,\dd E_A(\lambda),
 \qquad
 (z-A)^{-1}=\int_{\sigma(A)}\frac{\dd E_A(\lambda)}{z-\lambda}.}
 \label{eq:resolvent-spectral-resolution-dp68}
\end{equation}
由此可把预解式的算符范数直接写成到谱的距离：
\begin{equation}
\boxed{
\|(z-QHQ)^{-1}\|
=\frac1{\operatorname{dist}[z,\sigma(QHQ)]}}
\qquad
(z\notin\sigma(QHQ),\ QHQ=QHQ^\dagger).
\label{eq:resolvent-norm-spectral-dp21}
\end{equation}

若 $A$ 可逆，则
\begin{equation}
 \boxed{A-B=(I-BA^{-1})A.}
 \label{eq:neumann-factorization-dp68}
\end{equation}
当 $\|BA^{-1}\|<1$ 时，几何级数在算符范数下收敛，得到
\begin{equation}
\boxed{
(A-B)^{-1}=A^{-1}\sum_{n=0}^{\infty}(BA^{-1})^n}
\qquad
\text{if }\|BA^{-1}\|<1.
\label{eq:operator-neumann-resolvent-dp21}
\end{equation}

若 $P$ 中只保留一个非简并裸态 $|p\rangle$，且 $P$--$Q$ 混合足够弱，精确 Feshbach 自能的最低非零静态极限为
\begin{equation}
 \boxed{
 \Delta E_p^{(2)}
 =\sum_{q\in Q}
 \frac{|H_{pq}|^2}{E_p-E_q}.}
 \label{eq:feshbach-second-order-static-dp68}
\end{equation}
这就是普通非简并二阶微扰论在投影预解式中的来源。

\begin{derivation}[从\eqnrefs{eq:feshbach-block-matrix-dp68}到\eqnrefs{eq:feshbach-schur-proof-r10}]

正文\eqnrefs{eq:feshbach-block-matrix-dp68} 把完整预解式的逆算符分成目标块 $P$ 和补块 $Q$。只要 $z-QHQ$ 可逆，就能先从下方分块方程消去 $Q$ 分量，再代回 $P$ 方程；这一普通分块求逆步骤直接给出正文\eqnrefs{eq:feshbach-schur-proof-r10}，不需要弱耦合假设。

记
\begin{equation*}
 z-H=
 \begin{pmatrix}A&B\\C&D\end{pmatrix},
 \quad
 A=z-PHP,
 \quad B=-PHQ,
 \quad C=-QHP,
 \quad D=z-QHQ.
\end{equation*}
设其逆矩阵为
\begin{equation*}
 \begin{pmatrix}X&Y\\Z&W\end{pmatrix}.
\end{equation*}
由下左块方程
\begin{equation*}
 CX+DZ=0
\end{equation*}
得到
\begin{equation*}
 Z=-D^{-1}CX.
\end{equation*}
代入上左块 $AX+BZ=I_P$，
\begin{align*}
 (A-BD^{-1}C)X&=I_P,\\
 X&=(A-BD^{-1}C)^{-1}.
\end{align*}
而 $X=P(z-H)^{-1}P$。恢复四个分块以后即得正文\eqnrefs{eq:feshbach-schur-proof-r10}。
\end{derivation}

\begin{derivation}[从\eqnrefs{eq:sw-sylvester-r15}到\eqnrefs{eq:sw_arbitrary_target_subspace_matrix_elements}]
Baker--Campbell--Hausdorff 展开给出

\begin{align*}
 e^SHe^{-S}
 &=H+[S,H]+\frac12[S,[S,H]]+\order(S^3H)\\
 &=H_{\rm d}+H_{\rm od}+[S,H_{\rm d}]+[S,H_{\rm od}]
 +\frac12[S,[S,H_{\rm d}]]+\order(S^3H).
\end{align*}
由正文\eqnrefs{eq:sw-sylvester-r15}，$[S,H_{\rm d}]=-H_{\rm od}$，因此
\begin{align*}
 e^SHe^{-S}
 &=H_{\rm d}+[S,H_{\rm od}]
 -\frac12[S,H_{\rm od}]+\order(S^3H)\\
 &=H_{\rm d}+\frac12[S,H_{\rm od}]+\order(S^3H),
\end{align*}
即正文\eqnrefs{eq:sw-bch-second-r15}。

令
\begin{equation*}
 H_{\rm d}|p\rangle=E_p|p\rangle,
 \qquad
 H_{\rm d}|q\rangle=E_q|q\rangle,
\end{equation*}
其中 $p,p'\in P$、$q\in Q$。在 $\langle q|\cdots|p\rangle$ 上取正文\eqnrefs{eq:sw-sylvester-r15} 的矩阵元，
\begin{align*}
 0&=H_{qp}+(E_p-E_q)S_{qp},\\
 S_{qp}&=\frac{H_{qp}}{E_q-E_p},
 \qquad
 S_{pq}=\frac{H_{pq}}{E_p-E_q}.
\end{align*}
于是
\begin{align*}
 \langle p|[S,H_{\rm od}]|p'\rangle
 &=\sum_{q\in Q}
 \left(S_{pq}H_{qp'}-H_{pq}S_{qp'}\right)\\
 &=\sum_{q\in Q}H_{pq}H_{qp'}
 \left[
 \frac1{E_p-E_q}
 +\frac1{E_{p'}-E_q}
 \right].
\end{align*}
乘上正文\eqnrefs{eq:sw-bch-second-r15} 中的 $1/2$，并加回 $P H_{\rm d}P$，得到正文\eqnrefs{eq:sw_arbitrary_target_subspace_matrix_elements}。
\end{derivation}

\begin{derivation}[从\eqnrefs{eq:resolvent-spectral-resolution-dp68}到\eqnrefs{eq:resolvent-norm-spectral-dp21}]

正文\eqnrefs{eq:resolvent-spectral-resolution-dp68} 把自伴算符的预解式写成谱测度积分。对任意归一化态取范数平方，可以把算符范数问题化成正概率测度下的标量函数 $|z-\lambda|^{-2}$ 的上确界，从而得到正文\eqnrefs{eq:resolvent-norm-spectral-dp21}。

对归一化 $|\psi\rangle$ 定义
\begin{equation*}
 \dd\mu_\psi(\lambda)
 =\langle\psi|\dd E_A(\lambda)|\psi\rangle,
 \qquad
 \int\dd\mu_\psi=1.
\end{equation*}
于是
\begin{align*}
 \|(z-A)^{-1}|\psi\rangle\|^2
 &=\int_{\sigma(A)}
 \frac{\dd\mu_\psi(\lambda)}{|z-\lambda|^2}\\
 &\le
 \sup_{\lambda\in\sigma(A)}\frac1{|z-\lambda|^2}.
\end{align*}
对所有归一化态取上确界，得到
\begin{equation*}
 \|(z-A)^{-1}\|
 \le\frac1{\operatorname{dist}[z,\sigma(A)]}.
\end{equation*}
反向不等式可取一列谱权重越来越集中在离 $z$ 最近谱邻域的归一化态；离散谱情形就是取最近本征态。因此等号成立。令 $A=QHQ$ 即得正文\eqnrefs{eq:resolvent-norm-spectral-dp21}。
\end{derivation}

\begin{derivation}[从\eqnrefs{eq:feshbach_schur_exact_positive_projection}到\eqnrefs{eq:feshbach-second-order-static-dp68}]

精确 Feshbach--Schur 有效算符为
\[
H_{\rm eff}(z)
=
PHP+\Sigma_P(z),
\qquad
\Sigma_P(z)
=
PHQ\,(z-QHQ)^{-1}QHP.
\]
这里 \(P\) 可以是任意有限维目标子空间，并不要求只保留一个本征态。设目标能窗以参考能量 \(E_*\) 为中心，并与 \(QHQ\) 的谱保持最小距离
\[
\Delta_Q
\equiv
\operatorname{dist}\!\left(E_*,\sigma(QHQ)\right)>0.
\]
若实际极点满足 \(|z-E_*|<\Delta_Q\)，则
\[
z-QHQ
=
(E_*-QHQ)
\left[
I+(z-E_*)(E_*-QHQ)^{-1}
\right].
\]
定义
\[
X\equiv
(z-E_*)(E_*-QHQ)^{-1}.
\]
由预解式范数
\[
\|(E_*-QHQ)^{-1}\|
\le\frac1{\Delta_Q}
\]
可知
\[
\|X\|
\le\frac{|z-E_*|}{\Delta_Q}<1,
\]
所以在目标能窗内可以作算符 Neumann 展开：
\begin{align*}
(z-QHQ)^{-1}
&=
\left(I+X\right)^{-1}(E_*-QHQ)^{-1}\\
&=
\left(I-X+X^2-\cdots\right)
(E_*-QHQ)^{-1}\\
&=
(E_*-QHQ)^{-1}
-(z-E_*)(E_*-QHQ)^{-2}
+\mathcal O\!\left(
\frac{|z-E_*|^2}{\Delta_Q^3}
\right).
\end{align*}
因此
\begin{align*}
\Sigma_P(z)
&=
PHQ(E_*-QHQ)^{-1}QHP\\
&\quad
-(z-E_*)
PHQ(E_*-QHQ)^{-2}QHP
+\cdots .
\end{align*}

若 \(V\equiv PHQ+QHP\) 是弱混合，并定义
\[
\epsilon_{\rm mix}
\equiv
\frac{\|PHQ\|}{\Delta_Q}\ll1,
\]
则第一项满足
\[
\|\Sigma_P(E_*)\|
\le
\frac{\|PHQ\|^2}{\Delta_Q}
=
\epsilon_{\rm mix}^2\Delta_Q.
\]
精确极点相对未耦合目标能量的移动于是也是二阶，
\[
|z-E_*|
=
\mathcal O(\epsilon_{\rm mix}^2\Delta_Q).
\]
代回第二项可得
\[
\left\|
(z-E_*)
PHQ(E_*-QHQ)^{-2}QHP
\right\|
=
\mathcal O(\epsilon_{\rm mix}^4\Delta_Q),
\]
所以在二阶精度下，自能中的能量依赖可以冻结在 \(E_*\)：
\[
H_{\rm eff}^{(2)}
=
PHP
+
PHQ(E_*-QHQ)^{-1}QHP
+
\mathcal O(\epsilon_{\rm mix}^4\Delta_Q).
\]

正文\eqnrefs{eq:feshbach-second-order-static-dp68} 是一维目标子空间的特例。取
\[
P=|p\rangle\langle p|,
\qquad
PHP=E_pP,
\qquad
E_*=E_p,
\]
并对 \(QHQ\) 插入本征态完备关系
\[
Q
=
\sum_{q\in Q}|q\rangle\langle q|,
\qquad
QHQ|q\rangle=E_q|q\rangle,
\]
得到
\begin{align*}
\langle p|\Sigma_P(E_p)|p\rangle
&=
\sum_{q\in Q}
\frac{
\langle p|H|q\rangle
\langle q|H|p\rangle
}{E_p-E_q}\\
&=
\sum_{q\in Q}
\frac{|H_{pq}|^2}{E_p-E_q}.
\end{align*}
于是
\[
z
=
E_p
+
\sum_{q\in Q}
\frac{|H_{pq}|^2}{E_p-E_q}
+
\mathcal O(\epsilon_{\rm mix}^4\Delta_Q),
\]
即正文\eqnrefs{eq:feshbach-second-order-static-dp68}。因此“二阶能量修正”不是从非简并微扰论类比得到，而是精确 Schur 补在谱隙与弱混合控制下的静态展开；若目标子空间简并或近简并，应保留整个 \(P\) 块而不能先选单个 \(|p\rangle\)。
\end{derivation}


\section{函数积分与规范结构}

Hilbert 空间中的投影把自由度分成保留部分与消去部分；场论还需要处理自由度数目趋于无穷的极限。最稳妥的入口仍是有限维积分：先研究普通交换变量和反交换 Grassmann 变量的二次型积分，再把格点或模截断后的场看成一个高维向量。这样函数积分中的逆核、行列式和费米负号都先由有限维代数获得，再在连续极限中继承。连续对称与规范结构随后使用同样的“局部线性化”思想建立。

\subsection{高斯积分与反交换变量}
场论中的自由理论之所以特别重要，是因为它的作用量对场变量是二次的；二次型的逆给传播核，而带源二次积分又是后面 Wick 配对和微扰展开的母计算。玻色自由度使用普通交换变量即可；费米自由度若仍用普通数积分，就无法产生反交换统计和费米行列式，因此需要 Grassmann 代数与 Berezin 积分。先在有限维里把这两类二次积分算清楚，连续场函数积分只是在调节后把矩阵指标扩展为时空或模指标。

在有限维哈密顿系统中，正则对 $(q_a,p_a)$ 的量子化约定是
\begin{equation}
[\hat q_a,\hat p_b]=\ii\hbar\delta_{ab}.
\label{eq:global-canonical-quantization}
\end{equation}
场的正则结构不由“把离散指标换成连续坐标”推出。对玻色场，经典场论的辛形式直接给出等时 Poisson 括号
\[
 \{\phi_a(\mathbf x,t),\pi_b(\mathbf y,t)\}_{\rm PB}
 =\delta_{ab}\delta^{(3)}(\mathbf x-\mathbf y),
\]
正则量子化才把它变成 $[\hat\phi_a(\mathbf x,t),\hat\pi_b(\mathbf y,t)]=\ii\hbar\delta_{ab}\delta^{(3)}(\mathbf x-\mathbf y)$。费米场的作用量对时间导数是一阶的，约束正则结构与正能 Hamiltonian 要求使用反对易代数；费米场的反对易代数必须由其一阶正则结构和正能谱要求建立，不能由玻色对易子作形式替换得到。

先固定二次积分最常用的有限维恒等式。所谓 Grassmann 变量不是普通数，而是满足 $\theta_i\theta_j=-\theta_j\theta_i$、因而 $\theta_i^2=0$ 的代数变量；与它配套的 Berezin 积分本质上抽取变量的最高次系数。这些定义随后固定行列式与反交换符号规则。这里的有限维公式不能仅靠“把矩阵指标换成连续坐标”就自动成为无限维结论。对场函数作积分时，先用格点、模截断或其它有限自由度调节，使二次核成为有限维矩阵，再在调节极限中定义函数空间测度。有限维公式的作用，是展示调节后必然出现的代数结构：普通变量的二次核由逆矩阵控制二点响应，Grassmann 变量则额外产生行列式和反交换符号。量子场论的带源函数积分则把这一有限维结构应用到经过调节的物理作用量。

% Finite-dimensional Gaussian identities used before the path integral.

高斯积分是二次作用量与自由传播核的有限维原型。对普通实变量 $x$ 和 $a>0$，定义
\begin{equation}
 I(a)\equiv\int_{-\infty}^{\infty}\dd x\,\ee^{-ax^2}.
 \label{eq:real-gaussian-definition-dp68}
\end{equation}
这个积分的精确归一化为
\begin{equation}
 \boxed{I(a)=\sqrt{\frac{\pi}{a}}.}
 \label{eq:real-gaussian-r15}
\end{equation}
量子传播与 Fourier 型积分还会遇到纯振荡的二次相位。其积分值由加入正收敛因子后的绝对收敛积分连续延拓定义。在这一规定下，$A>0$ 时
\begin{equation}
 \boxed{
 \int\frac{\dd p}{2\pi\hbar}\,
 \ee^{-\ii A p^2/(2\hbar)}
 =\frac{1}{\sqrt{2\pi\ii\hbar A}}.}
 \label{eq:fresnel-normalized-r15}
\end{equation}
平方根的支由同一个收敛处方固定；传播核中的相位与归一化必须继承这一支选择。

% DeepPass68: finite-dimensional Gaussian and Grassmann identities.

多变量高斯公式直接由实对称正定二次型定义；一维积分只是其最简单的分量检验。对实对称正定矩阵 $A$，定义带源积分
\begin{equation}
\mathcal I_A[J]\equiv
\int_{\mathbb R^n}\dd^nq\,
\exp\!\left[-\frac12q^TAq+J^Tq\right].
\label{eq:euclidean-gaussian-definition-dp68}
\end{equation}
完成平方后得到
\begin{equation}
\boxed{
\mathcal I_A[J]
=(2\pi)^{n/2}(\det A)^{-1/2}
\exp\!\left[\frac12J^TA^{-1}J\right]}
\label{eq:euclidean-gaussian-source-dp21}
\end{equation}
以及归一化二点平均
\begin{equation}
\boxed{
\langle q_iq_j\rangle_A=(A^{-1})_{ij}.}
\label{eq:gaussian-covariance-dp21}
\end{equation}
这里 $\langle\cdots\rangle_A$ 表示以 $\ee^{-q^TAq/2}$ 为权重并除以零源积分后的平均。\eqnrefs{eq:gaussian-covariance-dp21} 体现了二次理论的基本结构：二次核 $A$ 控制作用，而二点响应由它的逆矩阵 $A^{-1}$ 控制。

费米自由度需要另一种积分。Grassmann 变量满足 $\theta_i\theta_j=-\theta_j\theta_i$，所以特别有 $\theta_i^2=0$。Berezin 积分不是面积积分，而是抽取最高次 Grassmann 系数的线性运算；一维约定固定为
\begin{equation}
\int\dd\theta\,1=0,
\qquad
\int\dd\theta\,\theta=1.
\label{eq:berezin-definition-dp68}
\end{equation}
对 $n$ 对独立变量 $\theta_i,\bar\theta_i$，固定测度次序后定义
\begin{equation}
\mathcal Z_G[\xi,\bar\xi]\equiv
\int\prod_{i=1}^{n}\dd\bar\theta_i\dd\theta_i\,
\exp\!\left[-\bar\theta M\theta+\bar\xi\theta+\bar\theta\xi\right],
\label{eq:grassmann-source-definition-dp68}
\end{equation}
其中 $M$ 可逆，而 $\xi,\bar\xi$ 也是 Grassmann 源。采用与\eqnrefs{eq:grassmann-source-definition-dp68} 一致的固定测度约定，零源积分为
\begin{equation}
\boxed{
\mathcal Z_G[0,0]=\det M.}
\label{eq:grassmann-gaussian-determinant-dp68}
\end{equation}
加入源以后，平移变量给出
\begin{equation}
\boxed{
\frac{\mathcal Z_G[\xi,\bar\xi]}{\mathcal Z_G[0,0]}
=\exp\!\left(\bar\xi M^{-1}\xi\right).}
\label{eq:grassmann-source-gaussian-dp68}
\end{equation}
因此，若对 $\bar\xi_i$ 取左导数、对 $\xi_j$ 取右导数，则零源二点收缩为
\begin{equation}
\boxed{
\left.
\frac{1}{\mathcal Z_G[0,0]}
\frac{\overrightarrow\partial\mathcal Z_G}{\partial\bar\xi_i}
\frac{\overleftarrow\partial}{\partial\xi_j}
\right|_{\xi=\bar\xi=0}
=(M^{-1})_{ij}.}
\label{eq:grassmann-inverse-kernel-dp68}
\end{equation}
这正是费米子自由理论中“二次算符的逆成为传播核”的有限维原型；闭合费米子圈的额外负号则来自 Wick 收缩中的费米置换宇称，与高斯积分行列式是不同的代数来源。

最后还需要知道 Grassmann 测度怎样随线性变量变换改变。对
\begin{equation}
\theta_i'=A_{ij}\theta_j,
\label{eq:berezin-linear-transform-dp68}
\end{equation}
固定同一测度次序后有
\begin{equation}
\boxed{
\dd\theta_n'\cdots\dd\theta_1'
=(\det A)^{-1}\,
\dd\theta_n\cdots\dd\theta_1.}
\label{eq:berezin-jacobian-dp21}
\end{equation}
它与普通变量的 Jacobian 方向相反，来源就是\eqnrefs{eq:berezin-definition-dp68} 的系数抽取规则。












\begin{derivation}[从\eqnrefs{eq:berezin-definition-dp68}到\eqnrefs{eq:grassmann-gaussian-determinant-dp68}]
Grassmann 高斯积分产生行列式的原因可以直接从“积分抽取最高次系数”看出。因为每个 $\theta_i$ 和 $\bar\theta_i$ 都平方为零，指数
\begin{equation*}
 \ee^{-\bar\theta M\theta}
 =\sum_{r=0}^{n}\frac{(-1)^r}{r!}(\bar\theta M\theta)^r
\end{equation*}
是有限多项式。测度含全部 $2n$ 个变量，所以积分后只有 $r=n$ 项可能留下。

展开 $(\bar\theta M\theta)^n$ 时，只要两个行指标或两个列指标重复，相应单项式就含某个 Grassmann 变量的平方而消失。因此非零项恰由 $1,\ldots,n$ 的排列标记。把变量重排回正文固定的测度次序会产生排列符号，而第一组指标的 $n!$ 种排列抵消展开中的 $1/n!$。于是
\begin{align*}
 \mathcal Z_G[0,0]
 &=\sum_{\sigma\in S_n}\operatorname{sgn}(\sigma)
 \prod_{i=1}^{n}M_{i\sigma(i)}\\
 &=\det M,
\end{align*}
最后一行正是行列式的 Leibniz 定义，从而得到正文\eqnrefs{eq:grassmann-gaussian-determinant-dp68}。
\end{derivation}

\begin{derivation}[从\eqnrefs{eq:grassmann-source-definition-dp68}到\eqnrefs{eq:grassmann-source-gaussian-dp68}]
加入源以后，Grassmann 高斯积分与普通高斯一样可以完成平方，只是变量和源的次序必须保持固定。令
\begin{equation*}
 \theta'=\theta-M^{-1}\xi,
 \qquad
 \bar\theta'=\bar\theta-\bar\xi M^{-1},
\end{equation*}
也就是 $\theta=\theta'+M^{-1}\xi$、$\bar\theta=\bar\theta'+\bar\xi M^{-1}$。代回指数并保持因子次序，有
\begin{align*}
 &-\bar\theta M\theta+\bar\xi\theta+\bar\theta\xi\\
 &=-\bar\theta'M\theta'
 -\bar\theta'\xi-\bar\xi\theta'-\bar\xi M^{-1}\xi
 +\bar\xi\theta'+\bar\xi M^{-1}\xi
 +\bar\theta'\xi+\bar\xi M^{-1}\xi\\
 &=-\bar\theta'M\theta'+\bar\xi M^{-1}\xi.
\end{align*}
Berezin 积分对这种常数 Grassmann 平移保持不变，因此
\begin{align*}
 \mathcal Z_G[\xi,\bar\xi]
 &=\ee^{\bar\xi M^{-1}\xi}
 \int\prod_i\dd\bar\theta_i'\dd\theta_i'\,
 \ee^{-\bar\theta'M\theta'}\\
 &=\mathcal Z_G[0,0]\ee^{\bar\xi M^{-1}\xi},
\end{align*}
即正文\eqnrefs{eq:grassmann-source-gaussian-dp68}。
\end{derivation}

\begin{derivation}[从\eqnrefs{eq:grassmann-source-gaussian-dp68}到\eqnrefs{eq:grassmann-inverse-kernel-dp68}]
正文已经固定了源的次序，因此二点收缩只需按同一约定求两次 Grassmann 导数。先对 $\bar\xi_i$ 取左导数，
\begin{equation*}
 \frac{\overrightarrow\partial\mathcal Z_G}{\partial\bar\xi_i}
 =\mathcal Z_G(M^{-1})_{ik}\xi_k.
\end{equation*}
再从右边对 $\xi_j$ 求导，并令源为零，指数因子回到 $1$，从而
\begin{equation*}
 \left.
 \frac{1}{\mathcal Z_G[0,0]}
 \frac{\overrightarrow\partial\mathcal Z_G}{\partial\bar\xi_i}
 \frac{\overleftarrow\partial}{\partial\xi_j}
 \right|_{0}
 =(M^{-1})_{ij}.
\end{equation*}
这就是正文\eqnrefs{eq:grassmann-inverse-kernel-dp68}。若交换源或导数的次序，会因 Grassmann 奇偶性出现额外负号，所以以后把这种带 Grassmann 源的高斯积分用于费米场时，必须同时固定源、导数和测度的次序。
\end{derivation}

\begin{derivation}[从\eqnrefs{eq:berezin-linear-transform-dp68}到\eqnrefs{eq:berezin-jacobian-dp21}]
普通变量的 Jacobian 来自体积元，而 Berezin 测度由最高次系数的归一化决定，所以变换方向恰好相反。先看一维 $\theta'=a\theta$。若写 $\dd\theta'=C\dd\theta$，则\eqnrefs{eq:berezin-definition-dp68} 要求
\begin{equation*}
 1=\int\dd\theta'\,\theta'
 =Ca\int\dd\theta\,\theta=Ca,
\end{equation*}
故 $C=a^{-1}$。

多维时由\eqnrefs{eq:berezin-linear-transform-dp68} 有
\begin{equation*}
 \theta_1'\cdots\theta_n'
 =A_{1i_1}\cdots A_{ni_n}\theta_{i_1}\cdots\theta_{i_n}.
\end{equation*}
只有 $(i_1,\ldots,i_n)$ 为一个排列时该项非零；重排到固定次序所产生的符号正好组成行列式，因此
\begin{equation*}
 \theta_1'\cdots\theta_n'=(\det A)\theta_1\cdots\theta_n.
\end{equation*}
若 $\dd^n\theta'=C\dd^n\theta$，最高次系数的 Berezin 归一化给出 $C\det A=1$，于是得到正文\eqnrefs{eq:berezin-jacobian-dp21}。
\end{derivation}


\subsection{连续对称与规范几何}
高斯与反交换变量积分解决“怎样对许多自由度积分”，但它还没有回答另一类问题：若某个物理系统在一整族连续变换下保持不变，怎样系统记录这些变换的组合规律，以及同一个变换怎样作用在普通数值场、时空向量或内部多分量场上？先从“群”这一最基本的变换集合开始定义，再逐步引入连续群的无穷小生成方向、表示以及局域规范场所需的联络与曲率。所有术语第一次出现时都按普通线性代数和多元微积分可以理解的方式建立；$U(1)$、$SU(2)$、$SU(3)$ 和 Lorentz 对称随后只是这套一般语言的具体实例。

空间转动改变坐标分量却保持欧氏距离，量子态的整体相位变换改变复相位却保持全部概率；这两件熟悉的事实都具有同一结构：允许的变换可以连续复合，而且复合后物理不变量保持不变。群论把这种“变换怎样复合”抽象出来。Lorentz 群、$U(1)$、$SU(2)$ 和 $SU(3)$ 的差别在于各自保持什么对象不变，但它们都服从同一组群运算规则。

先设 $G$ 是一组可以依次执行的变换，并把“先做 $g_2$、再做 $g_1$”记成乘积 $g_1g_2$。若这组变换满足
\begin{equation}
 g_1g_2\in G,\qquad
 (g_1g_2)g_3=g_1(g_2g_3),\qquad
 eg=ge=g,\qquad
 gg^{-1}=g^{-1}g=e,
 \label{eq:group-axioms-math-dp76}
\end{equation}
就称 $G$ 为一个\emph{群}。四个条件依次表示：两个允许变换连续做仍是允许变换；连续做三次时括号位置不影响结果；存在“什么也不做”的单位变换 $e$；每个变换都存在可以把它撤销的逆变换。群不要求乘法可交换，所以一般可以有 $g_1g_2\neq g_2g_1$。例如三维空间绕不同轴的有限转动通常不交换，这已经是非交换群最直观的例子。

当群元素可以由若干连续实参数平滑标记时，离散的复合规则还能进一步被微分。严格地说，若群元素构成光滑流形，而且群乘法与取逆都是光滑映射，就称为\emph{Lie 群}。这里“流形”只需理解成：每个群元素附近都能选取若干实数作为局部坐标，并对穿过该点的群曲线像普通多元函数那样求导。$SO(1,3)$、$U(1)$ 和 $SU(N)$ 都可表示为满足光滑矩阵约束的有限维矩阵集合，因此其无穷小变换可以直接用矩阵微分处理。

Lie 群最重要的简化发生在单位元附近。取一条通过单位元的光滑群曲线 $g(t)$，满足 $g(0)=e$。它在 $t=0$ 的速度
\begin{equation*}
 X\equiv\left.\frac{\dd g(t)}{\dd t}\right|_{t=0}
\end{equation*}
描述“沿某个方向做无穷小变换”。所有这类速度组成一个实向量空间，称为单位元处的\emph{切空间}：
\begin{equation}
 \boxed{\mathfrak g\equiv T_eG.}
 \label{eq:lie-algebra-tangent-general-dp68}
\end{equation}
这个向量空间连同其 Lie 括号称为群 $G$ 的\emph{Lie 代数}。它不是另一套独立对称性，而是有限群变换在单位元附近的一阶微分结构。选择一组基 $\{X_a\}_{a=1}^{d_G}$，就等于选择 $d_G$ 个独立的无穷小变换方向，其中
\begin{equation}
 d_G\equiv\dim_{\mathbb R}\mathfrak g
 \label{eq:lie-algebra-dimension-dp68}
\end{equation}
是这些独立方向的数目。于是单位元附近的一般变换可写成
\begin{equation}
 g(\boldsymbol\theta)
 =I+\theta^aX_a+\order(\theta^2).
 \label{eq:lie-infinitesimal-dp11}
\end{equation}
这里的 $X_a$ 是群流形单位元处的切向量；当群进一步表示为量子 Hilbert 空间上的幺正算符时，相应切向表示才给出熟悉的厄米量子生成元。

若只沿一个固定方向 $X$ 连续变换，就得到单参数子群 $g(t)$，它满足
\begin{equation}
 g(t_1)g(t_2)=g(t_1+t_2),
 \qquad g(0)=I.
 \label{eq:one-param-group-dp11}
\end{equation}
对矩阵 Lie 群，这个函数由初始速度 $X=\dot g(0)$ 唯一确定为
\begin{equation}
 \boxed{g(t)=\exp(tX),}
 \label{eq:one-param-exp-dp11}
\end{equation}
其中矩阵指数就是普通幂级数
\begin{equation}
 e^X\equiv\sum_{n=0}^{\infty}\frac{X^n}{n!}.
 \label{eq:matrix-exponential-series-dp68}
\end{equation}
因此“生成元”这个词的含义很具体：知道单位元处的无穷小方向 $X$，指数映射就沿这个方向生成有限连续变换。

连续群还必须记录不同无穷小方向的先后次序是否重要。对所用的矩阵 Lie 群，群的共轭映射会把单位元切空间映回自身，因此两个生成方向的矩阵对易子仍是合法的无穷小生成方向。矩阵 Lie 群的 Lie 括号于是为
\begin{equation}
 \boxed{[X,Y]_{\mathfrak g}=XY-YX,\qquad X,Y\in\mathfrak g.}
 \label{eq:matrix-lie-bracket-math-dp76}
\end{equation}
它测量两个无穷小变换方向的非交换部分；若该括号恒为零，相应 Lie 代数就是阿贝尔的。

在基 $X_a$ 上展开后，必存在一组实数 $f_{ab}{}^c$ 使
\begin{equation}
 \boxed{[X_a,X_b]_{\mathfrak g}=f_{ab}{}^cX_c.}
 \label{eq:lie-algebra-general-dp11}
\end{equation}
这些数称为\emph{结构常数}：它们记录“先沿第 $a$ 个无穷小方向、再沿第 $b$ 个方向”所产生的非交换效应如何重新分解到原来的生成方向上。矩阵乘法的结合律使对易子满足 Jacobi 恒等式
\begin{equation}
 [A,[B,C]]+[B,[C,A]]+[C,[A,B]]=0,
 \label{eq:commutator-jacobi-dp68}
\end{equation}
因而
\begin{equation}
 \boxed{
 f_{bc}{}^{d}f_{ad}{}^{e}
 +f_{ca}{}^{d}f_{bd}{}^{e}
 +f_{ab}{}^{d}f_{cd}{}^{e}=0.}
 \label{eq:jacobi-structure-dp11}
\end{equation}
Baker--Campbell--Hausdorff 公式只是把同一件事提升到有限但仍靠近单位元的变换：
\begin{equation}
 e^Xe^Y
 =\exp\!\left[
 X+Y+\frac12[X,Y]
 +\frac1{12}[X,[X,Y]]+\frac1{12}[Y,[Y,X]]+\cdots
 \right].
 \label{eq:bch-lie-dp11}
\end{equation}
它说明有限变换的非交换修正由同一个 Lie 括号逐级组织。

抽象群只规定变换怎样彼此复合；物理计算还需要知道同一个变换怎样作用在具体对象上。例如空间转动既可以作用在三维位置向量上，也可以作用在某个具有多个复分量的量子态向量上；两者使用的矩阵维数可以不同，却描述同一个抽象转动。\emph{表示}就是把这种“同一抽象变换在某个向量空间上的具体作用”数学化。令 $V_R$ 是一组物理分量所在的复向量空间，$GL(V_R)$ 表示 $V_R$ 上所有可逆线性变换组成的群。一个群表示是映射
\begin{equation*}
 D_R:G\longrightarrow GL(V_R)
\end{equation*}
并满足
\begin{equation*}
 D_R(g_1g_2)=D_R(g_1)D_R(g_2),
 \qquad D_R(e)=I_R.
\end{equation*}
第一条条件保证“先后做两个抽象对称变换”与“把对应矩阵先后作用在物理分量上”完全一致。数学上，把乘法关系保持不变的映射称为\emph{群同态}；因此群表示就是一个取值在可逆线性变换群中的群同态。表示空间的复维数记为
\begin{equation*}
 d_R\equiv\dim_{\mathbb C}V_R.
\end{equation*}
如果存在一个既非零也非整个 $V_R$ 的子空间，并且所有 $D_R(g)$ 都把它映回自身，那么表示还能继续分块；若不存在这样的真不变子空间，就称该表示为\emph{不可约表示}。不可约表示是对称性能够区分的最小独立块，所以后面“粒子按质量、自旋分类”最终会落到寻找 Poincar\'e 群的不可约表示。

量子内部对称通常要求保持态的内积，因此采用酉表示。此时表示在单位元附近的导数 $dD_R(X_a)$ 是反厄米算符。为了使用量子力学熟悉的厄米可观测量 convention，定义
\begin{equation}
 \boxed{T_R^a\equiv \ii\,dD_R(X_a),}
 \qquad
 D_R(\ee^{\theta^aX_a})
 =\exp(-\ii\theta^aT_R^a).
 \label{eq:unitary-representation-generators-dp68}
\end{equation}
因为表示保持群乘法，它的微分也保持 Lie 括号：
\begin{equation*}
 dD_R([X_a,X_b])
 =[dD_R(X_a),dD_R(X_b)].
\end{equation*}
于是具体表示中的厄米生成元满足
\begin{equation}
 \boxed{[T_R^a,T_R^b]=\ii f_{ab}{}^{c}T_R^c.}
 \label{eq:rep-lie-algebra-dp11}
\end{equation}
这里要区分三个层级：$X_a$ 是抽象群的无穷小方向，$f_{ab}{}^c$ 描述群本身的局部乘法结构，而 $T_R^a$ 是这些方向在特定表示空间 $V_R$ 上的矩阵实现。改变表示会改变矩阵 $T_R^a$，但不会改变抽象群是哪一个群。



生成元的对易关系描述群的局部乘法结构，但不能单独区分不同表示的尺度与二次不变量。为此需要基底无关的表示不变量。群若作为参数空间是紧致的，直观上表示其连续参数不会沿某个方向无限逃逸；$U(N)$ 与 $SU(N)$ 是典型紧致矩阵群。Lie 代数若除了 $\{0\}$ 和整个代数外，不存在对所有外部生成元作 Lie 括号后仍封闭的非零真子空间，就称为简单 Lie 代数；这样的封闭子空间称为理想。$SU(2)$ 与 $SU(3)$ 的 Lie 代数都是紧致简单 Lie 代数，因此其不可约表示可以用一组标准二次不变量比较。

在紧致简单 Lie 代数上，可以选择一个被全部群变换保持的不变内积，并把基归一化到该内积的分量为 $\delta^{ab}$。这相当于为“生成元方向之间怎样取长度和夹角”固定统一标尺。在不可约酉表示 $R$ 中，生成元的二次迹只能与这个不变内积成正比：
\begin{equation}
 \boxed{
 \Tr_R(T_R^aT_R^b)=T(R)\delta^{ab}.}
 \label{eq:dynkin-index-dp11}
\end{equation}
比例系数 $T(R)$ 衡量表示中生成元的整体迹归一化；它通常称为表示 $R$ 的 Dynkin 指标。它不是脱离生成元归一化的绝对数值：若所有生成元整体重标度，$T(R)$ 与二次 Casimir $C_2(R)$ 会同步改变。

同一组生成元还定义算符
\begin{equation}
 \mathcal C_R\equiv\sum_{a=1}^{d_G}T_R^aT_R^a.
 \label{eq:quadratic-invariant-operator-dp68}
\end{equation}
利用结构常数对生成元指标的反对称性可得 $[\mathcal C_R,T_R^b]=0$。这里使用一个基本线性代数事实：若一个复不可约表示中的线性算符与所有表示矩阵都对易，那么它不能在表示空间内再挑出新的不变方向，只能是恒等算符的常数倍。这就是 Schur 引理。因此
\begin{equation}
 \boxed{
 \mathcal C_R=C_2(R)I_{d_R}.}
 \label{eq:casimir-general-rep-dp11}
\end{equation}
这个与全部生成元对易的二次算符通常称为二次 Casimir 算符，$C_2(R)$ 是它在表示 $R$ 上的本征值。对可约表示，应在每个不可约块上分别使用这一式。

现在对\eqnrefs{eq:casimir-general-rep-dp11} 取整个表示空间的迹，并用\eqnrefs{eq:dynkin-index-dp11}：
\begin{align}
 d_RC_2(R)
 &=\sum_{a=1}^{d_G}\Tr_R(T_R^aT_R^a),\\
 &=T(R)\sum_{a=1}^{d_G}\delta^{aa}
 =d_GT(R).
\end{align}
于是得到一般恒等式
\begin{equation}
 \boxed{d_RC_2(R)=d_GT(R).}
 \label{eq:casimir-dynkin-relation-dp11}
\end{equation}
这里三个对象必须区分：$d_G=\dim_{\mathbb R}\mathfrak g$ 是该简单 Lie 代数的生成元数目，$d_R=\dim_{\mathbb C}V_R$ 是表示空间维数，而 $T(R)$ 与 $C_2(R)$ 是在既定生成元归一化下定义的表示不变量。\eqnrefs{eq:casimir-dynkin-relation-dp11} 是后面所有 $C_F,C_A,T_F$ 数值的母关系，而不是由这些特殊数值反推的一般化。

最简单的紧 Lie 群 $U(1)$ 只有一个生成方向，因此 Lie 代数是阿贝尔的。由于阿贝尔生成元可以与耦合常数作相反的整体重标度，$U(1)$ 本身没有由非阿贝尔结构固定的唯一归一化。在一维不可约酉表示中，生成元只是实数 $y$，可写成
\begin{equation}
 \psi\longmapsto \ee^{-\ii y\theta}\psi.
 \label{eq:u1-charge-rep-dp11}
\end{equation}
所以阿贝尔“荷”在数学上就是选定生成元归一化后的一维表示标签。具体理论可以把这个实数解释成不同的守恒荷；它究竟对应哪种物理电荷，必须等作用量和相互作用建立后再由理论本身确定。

现在把抽象结构代入最常见的矩阵群。$U(N)$ 是所有 $N\times N$ 酉矩阵组成的群，即 $U^\dagger U=I_N$；其中再要求行列式为 $1$ 的子群称为 $SU(N)$，字母 $S$ 表示 ``special''。最小的非阿贝尔例子是 $SU(2)$：

$SU(2)$ 定义为
\begin{equation*}
 SU(2)=\{U\in\mathbb C^{2\times2}\mid U^\dagger U=I_2,\ \det U=1\}.
\end{equation*}
所谓\emph{基本表示}，就是让群矩阵直接作用在定义它的列向量空间 $\mathbb C^2$ 上；它的维数为 $2$。在这一本征的二维表示中，可取厄米生成元 $T^a=\sigma^a/2$。Pauli 矩阵满足
\begin{equation}
 \sigma^a\sigma^b
 =\delta^{ab}I_2+\ii\epsilon^{abc}\sigma^c,
 \label{eq:pauli-product-lie-dp68}
\end{equation}
因此
\begin{equation}
 \boxed{[T^a,T^b]=\ii\epsilon^{abc}T^c.}
 \label{eq:su2-algebra-dp11}
\end{equation}
定义 $T^\pm=T^1\pm\ii T^2$ 后，
\begin{equation}
 [T^3,T^\pm]=\pm T^\pm,
 \qquad [T^+,T^-]=2T^3.
 \label{eq:su2-ladder-algebra-dp11}
\end{equation}
一般不可约表示可由半整数 $j$ 标记，其二次 Casimir 与 $T^3$ 本征值为
\begin{equation}
 \boxed{
 C_2(j)=j(j+1),
 \qquad m=-j,-j+1,\ldots,j.}
 \label{eq:su2-casimir-dp11}
\end{equation}
$j=1/2$ 给二维基本表示，$j=1$ 给三维表示。这里的 $j$ 只标记抽象 $SU(2)$ 表示；两个具有同构 Lie 代数的 $SU(2)$ 对称完全可以作用在不同的向量空间上，因此表示标签本身还不指定具体物理含义。

同理，对一般 $SU(N)$，基本表示就是矩阵 $U$ 直接作用在 $\mathbb C^N$ 上，所以表示空间维数为 $d_F=N$。群还可以作用在自己的无穷小生成方向上：给定一个 Lie 代数元素 $X$，有限群元素 $U$ 通过共轭变换 $X\mapsto UXU^{-1}$ 把它线性混合成其它生成方向。因为被变换的向量空间就是 Lie 代数本身，这个表示称为\emph{伴随表示}。对 $SU(2)$，上面的 $j=1$ 三维表示正是这一伴随表示。一般 $SU(N)$ 的 Lie 代数有 $d_G$ 个独立生成方向，因此伴随表示的维数也是
\begin{equation*}
 d_G=N^2-1.
\end{equation*}
粒子物理通常把基本表示生成元的归一化约定固定为
\begin{equation}
 \Tr_F(T^aT^b)=T_F\delta^{ab},
 \qquad \boxed{T_F\equiv T(F)=\frac12.}
 \label{eq:generator-trace-normalization-dp11}
\end{equation}
这里 $T_F=1/2$ 是生成元归一化的选择，而不是先于一般表示理论的独立动力学结论。代入一般恒等\eqnrefs{eq:casimir-dynkin-relation-dp11} 立即得到基本表示的二次 Casimir
\begin{equation}
 \boxed{C_F\equiv C_2(F)=\frac{N^2-1}{2N}.}
 \label{eq:cf-general-math-dp11}
\end{equation}
同一归一化下，伴随表示生成元可写成 $(F^a)_{bc}=-\ii f^{abc}$；$SU(N)$ 的伴随二次 Casimir 为
\begin{equation}
 \boxed{C_A\equiv C_2(\mathrm{Adj})=N.}
 \label{eq:ca-general-math-dp11}
\end{equation}
基本生成元还满足完备关系
\begin{equation}
 \boxed{
 \sum_{a=1}^{N^2-1}(T^a)_{ij}(T^a)_{kl}
 =T_F\left(\delta_{il}\delta_{jk}-\frac1N\delta_{ij}\delta_{kl}\right).}
 \label{eq:suN-fierz-math-dp68}
\end{equation}
它只属于 $SU(N)$ 基本表示这一特例，用于把基本表示生成元的求和化为 Kronecker 张量。作为一个常用的 $N=3$ 特例，有
\begin{equation*}
 T_F=\frac12,\qquad C_F=\frac43,\qquad C_A=3.
\end{equation*}
因此这些数值只是一般 $SU(N)$ 表示结构在 $N=3$ 的基本表示与伴随表示上的末端代入。

相互作用项同时含多个场，因此还需要张量积表示。若 $v\in V_R$、$w\in V_S$，则
\begin{equation}
 D_{R\otimes S}(g)=D_R(g)\otimes D_S(g),
 \label{eq:tensor-product-rep-dp11}
\end{equation}
其无穷小生成元为
\begin{equation}
 \boxed{
 T_{R\otimes S}^a
 =T_R^a\otimes I+I\otimes T_S^a.}
 \label{eq:tensor-product-generator-dp11}
\end{equation}
若张量积包含一维平凡表示，就存在不变张量把各表示指标收缩成群不变量。局域多场乘积能够进入对称作用量的条件，正是其表示指标能够借助这类不变张量收缩成单态。

对任意 $2\times2$ 矩阵 $U$，完全反对称张量
\begin{equation*}
 \epsilon=\begin{pmatrix}0&1\\-1&0\end{pmatrix}=\ii\sigma^2
\end{equation*}
满足
\begin{equation}
 U^T\epsilon U=(\det U)\epsilon.
 \label{eq:2x2-epsilon-det-dp68}
\end{equation}
因此对 $U\in SU(2)$，
\begin{equation}
 \boxed{U^T\epsilon U=\epsilon.}
 \label{eq:su2-epsilon-invariance-dp11}
\end{equation}
这意味着 $u^T\epsilon v$ 是 $SU(2)$ 单态。共轭二重态定义为
\begin{equation}
 \boxed{\widetilde\Phi\equiv\ii\sigma^2\Phi^*,}
 \label{eq:tilde-higgs-math-dp11}
\end{equation}
并仍按基本二重态变换。这个构造只依赖 $SU(2)$ 的不变反对称张量，后面遇到需要把共轭场重新组织成二重态的物理模型时可以直接使用。

对 $SU(N)$ 的基本表示与反基本表示，$\delta_i{}^j$ 是不变张量。三维完全反对称张量与行列式满足
\begin{equation}
 U_i{}^{i'}U_j{}^{j'}U_k{}^{k'}\epsilon_{i'j'k'}
 =(\det U)\epsilon_{ijk}.
 \label{eq:determinant-epsilon-identity-dp68}
\end{equation}
因此对 $SU(3)$，
\begin{equation}
 \boxed{
 U_i{}^{i'}U_j{}^{j'}U_k{}^{k'}\epsilon_{i'j'k'}
 =\epsilon_{ijk}.}
 \label{eq:su3-epsilon-invariant-dp11}
\end{equation}
三个基本三维表示因而可以通过 $\epsilon_{ijk}$ 收缩成群单态。

除二次迹外，还可以从三个生成元构造完全对称的表示不变量
\begin{equation}
 d_R^{abc}
 \equiv2\Tr_R\!\left[T_R^a\{T_R^b,T_R^c\}\right].
 \label{eq:anomaly-symmetric-trace-dp11}
\end{equation}
它只由表示 $R$ 和生成元归一化决定。先把它作为群表示的三阶对称张量保存下来；它在量子理论中的具体用途等相应圈图与量子对称性问题建立后再解释。

Lie 群本身的局部结构由单位元切空间与 Lie 括号决定；表示把同一抽象 Lie 代数映到具体内部向量空间；$T(R)$ 与 $C_2(R)$ 是在既定生成元归一化下定义的二次表示不变量，并先满足 $d_RC_2(R)=d_GT(R)$，随后才有 $SU(N)$ 基本表示中的 $T_F,C_F,C_A$。张量积与不变张量决定多场局域项能否组成群单态。

Lie 群告诉我们在同一个时空点上允许怎样旋转内部向量，但局域规范变换允许这个内部基底随位置改变。于是还缺一个新的规则：怎样比较 $x$ 点与 $x+\dd x$ 点的内部向量。规范联络正是这条比较规则；曲率则回答“沿不同路径比较以后，结果是否一致”。在进入联络之前，先用微分形式把反对称指标和边界积分写成一种不依赖坐标排列的语言。

对 $p$ 阶微分形式
\begin{equation*}
 \omega^{(p)}
 =\frac1{p!}\omega_{\mu_1\cdots\mu_p}(x)
 \dd x^{\mu_1}\wedge\cdots\wedge\dd x^{\mu_p},
\end{equation*}
楔积满足 $\dd x^\mu\wedge\dd x^\nu=-\dd x^\nu\wedge\dd x^\mu$。外微分定义为
\begin{equation}
 \dd\omega^{(p)}
 =\frac1{p!}\partial_\nu\omega_{\mu_1\cdots\mu_p}
 \dd x^\nu\wedge\dd x^{\mu_1}\wedge\cdots\wedge\dd x^{\mu_p}.
 \label{eq:exterior-derivative-definition-dp68}
\end{equation}
由偏导交换与楔积反对称性得到两条基本规则：
\begin{equation}
 \boxed{
 \dd^2=0,
 \qquad
 \dd(\alpha^{(p)}\wedge\beta)
 =(\dd\alpha^{(p)})\wedge\beta
 +(-1)^p\alpha^{(p)}\wedge\dd\beta.}
 \label{eq:exterior-calculus-identities-dp68}
\end{equation}
与局域恒等式 $\dd^2=0$ 配套的全局积分关系是 Stokes 定理
\begin{equation}
 \boxed{\displaystyle \int_M\dd\omega=\int_{\partial M}\omega.}
 \label{eq:stokes-forms-dp11}
\end{equation}
它把“体积分变边界项”的计算统一为同一条数学定理。

现在回到局域内部空间。令 $T^a$ 为厄米生成元，矩阵值联络一形式与协变导数定义为
\begin{equation}
 \boxed{
 \mathscr C=\mathscr C_\mu^aT^a\dd x^\mu,
 \qquad
 D=\dd+\ii g_{\rm G}\mathscr C.}
 \label{eq:connection-oneform-dp11}
\end{equation}
$\mathscr C$ 的作用是规定相邻点的内部向量怎样比较。连续作用两次 $D$ 后，路径依赖的最低阶部分定义曲率：
\begin{equation}
 \boxed{
 \mathcal F\equiv\dd\mathscr C+\ii g_{\rm G}\mathscr C\wedge\mathscr C,
 \qquad
 D^2=\ii g_{\rm G}\mathcal F.}
 \label{eq:curvature-form-dp11}
\end{equation}
把二形式展开成分量，
\begin{equation}
 \boxed{
 \mathcal F_{\mu\nu}
 =\partial_\mu\mathscr C_\nu-\partial_\nu\mathscr C_\mu
 +\ii g_{\rm G}[\mathscr C_\mu,\mathscr C_\nu].}
 \label{eq:curvature-components-dp68}
\end{equation}
阿贝尔理论中最后的对易子消失；非阿贝尔场强的额外项直接来自矩阵不交换。

局域基变换 $\Psi'=U\Psi$ 不能改变“平行”的物理含义，因此要求
\begin{equation}
 D'\Psi'=U(D\Psi).
 \label{eq:covariant-derivative-covariance-dp68}
\end{equation}
这唯一固定联络的有限变换律
\begin{equation}
 \boxed{
 \mathscr C'
 =U\mathscr C U^{-1}
 +\frac{\ii}{g_{\rm G}}(\dd U)U^{-1}.}
 \label{eq:finite-gauge-connection-dp11}
\end{equation}
并进一步给出曲率的协变变换
\begin{equation}
 \boxed{\mathcal F'=U\mathcal F U^{-1}.}
 \label{eq:curvature-gauge-covariance-dp68}
\end{equation}
所以 $\Tr(\mathcal F_{\mu\nu}\mathcal F^{\mu\nu})$ 在局域基变换下保持不变。

曲率也可以直接由协变导数的对易子读出：
\begin{equation}
 [D_\mu,D_\nu]=\ii g_{\rm G}\mathcal F_{\mu\nu}.
 \label{eq:covariant-commutator-curvature-dp68}
\end{equation}
把它代入三个协变导数的 Jacobi 恒等式，就得到
\begin{equation}
 \boxed{
 D_\lambda\mathcal F_{\mu\nu}
 +D_\mu\mathcal F_{\nu\lambda}
 +D_\nu\mathcal F_{\lambda\mu}=0.}
 \label{eq:nonabelian-bianchi-dp11}
\end{equation}
这就是非阿贝尔 Bianchi 恒等式。

联络的几何意义最直接地体现在沿曲线输运。对曲线 $C:x^\mu(s)$，平行输运方程为
\begin{equation}
 \frac{\dd\Psi}{\dd s}
 +\ii g_{\rm G}\dot x^\mu(s)\mathscr C_\mu[x(s)]\Psi(s)=0.
 \label{eq:parallel-transport-ode-dp68}
\end{equation}
不同参数点上的矩阵一般不交换，所以有限解必须保留乘法次序：
\begin{equation}
 \boxed{
 U_C(s_f,s_i)
 =\mathcal P\exp\!\left[-\ii g_{\rm G}\int_C\mathscr C\right].}
 \label{eq:wilson-line-general-dp11}
\end{equation}
这个对象称为 Wilson 线。规范变换只旋转它的两个端点，
\begin{equation}
 \boxed{
 U_C' =U(x_f)U_CU^{-1}(x_i).}
 \label{eq:wilson-endpoint-covariance-dp68}
\end{equation}
因此对闭合曲线 $x_f=x_i$ 取迹后，端点旋转在迹中相消，得到规范不变的 Wilson 圈
\begin{equation}
 \boxed{
 W_R(C)=\frac1{d_R}\Tr_R\,
 \mathcal P\exp\!\left[-\ii g_{\rm G}\oint_C\mathscr C\right].}
 \label{eq:wilson-loop-general-dp11}
\end{equation}

曲率为什么可以理解成“局部的不闭合量”，可以直接看一个无穷小矩形。沿四条边依次做平行输运并保留面积阶，
\begin{equation}
 \boxed{
 U_\square
 =I-\ii g_{\rm G}\mathcal F_{\mu\nu}(x)
 \delta x^\mu\delta x^\nu
 +\order(\delta x^3).}
 \label{eq:small-loop-curvature-dp11}
\end{equation}
所以 $\mathcal F_{\mu\nu}$ 正是闭合小回路输运偏离单位矩阵的面积密度。

前面的讨论都是局域的。还需要一个与非 Abel 规范场全局结构直接相关的结果：在四维 Euclidean 时空中，对作用量有限并在无穷远趋于纯规范的 $SU(N)$ 场，可以定义
\begin{equation}
 Q_{\rm top}
 \equiv\frac{g_{\rm G}^2}{8\pi^2}\int\Tr(\mathcal F\wedge\mathcal F).
 \label{eq:pontryagin-form-definition-dp68}
\end{equation}
用 $\Tr(T^aT^b)=\delta^{ab}/2$ 展开分量后，
\begin{equation}
 \boxed{
 Q_{\rm top}
 =\frac{g_{\rm G}^2}{32\pi^2}\int\dd^4 x\,
 \mathcal F_{\mu\nu}^a{}^\star\!\mathcal F^{a\mu\nu}
 \in\mathbb Z.}
 \label{eq:pontryagin-number-dp11}
\end{equation}
这里从微分形式到分量公式只是代数换写；$Q_{\rm top}$ 为整数则还需要无穷远边界条件与整体拓扑定理。也就是说，“归一化系数是多少”可以逐行推导，而“为什么只能取整数”不是局域代数能够单独证明的结论。

Lie 群告诉我们在单个点上允许怎样改变内部基底；联络规定从一个点走到相邻点时怎样比较两个基底；曲率则测量沿两个不同方向依次比较所产生的差异。阿贝尔电磁理论的内部空间是一维相位，矩阵自动交换；非阿贝尔规范理论的内部空间是多维的，比较顺序因此进入场强本身。



\begin{derivation}[从\eqnrefs{eq:one-param-group-dp11}到\eqnrefs{eq:matrix-lie-bracket-math-dp76}]
先沿固定切向方向 $X=\dot g(0)$ 考察单参数子群。由群关系 $g(t+s)=g(t)g(s)$ 对 $s$ 求导并令 $s=0$，得到
\begin{align*}
 \dot g(t)
 &=g(t)\left.\frac{\dd g(s)}{\dd s}\right|_{s=0}\\
 &=g(t)X.
\end{align*}
交换 $t,s$ 的角色，同理由 $g(t+s)=g(s)g(t)$ 得 $\dot g(t)=Xg(t)$，所以 $X$ 与该单参数子群上的 $g(t)$ 对易。设
\begin{equation*}
 g(t)=\sum_{n=0}^{\infty}t^nG_n,
\end{equation*}
代入 $\dot g=Xg$ 并比较 $t^n$ 的系数，得到
\begin{equation*}
 (n+1)G_{n+1}=XG_n,\qquad G_0=I.
\end{equation*}
递推给出 $G_n=X^n/n!$，于是
\begin{equation*}
 g(t)=I+tX+\frac{t^2X^2}{2!}+\cdots=\ee^{tX},
\end{equation*}
即正文\eqnrefs{eq:one-param-exp-dp11}。

接着取另一个切向量 $Y\in\mathfrak g=T_eG$，并选一条群曲线 $h(t)$ 满足 $h(0)=I$、$\dot h(0)=Y$。由刚得到的 $\ee^{sX}\in G$ 以及群乘法和取逆的封闭性，
\begin{equation*}
 h_s(t)=\ee^{sX}h(t)\ee^{-sX}\in G.
\end{equation*}
因此它在 $t=0$ 的切向量仍属于单位元切空间：
\begin{equation*}
 Y(s)=\left.\frac{\dd h_s(t)}{\dd t}\right|_{t=0}
 =\ee^{sX}Ye^{-sX}\in\mathfrak g.
\end{equation*}
对 $s$ 再求导，
\begin{align*}
 \frac{\dd Y(s)}{\dd s}
 &=Xe^{sX}Ye^{-sX}-\ee^{sX}Ye^{-sX}X,\\
 \left.\frac{\dd Y(s)}{\dd s}\right|_{s=0}
 &=XY-YX.
\end{align*}
因为 $Y(s)$ 是向量空间 $\mathfrak g$ 内的光滑曲线，其导数仍属于 $\mathfrak g$，所以
\begin{equation*}
 XY-YX\in\mathfrak g,
\end{equation*}
从而得到正文\eqnrefs{eq:matrix-lie-bracket-math-dp76}。
\end{derivation}

\begin{derivation}[从\eqnrefs{eq:lie-algebra-general-dp11}到\eqnrefs{eq:rep-lie-algebra-dp11}]
需要验证的关键点是：群表示在单位元处求导以后，仍然保持 Lie 括号。取两个抽象生成方向 $X,Y\in\mathfrak g$。群中的共轭变换满足
\[
\operatorname{Ad}_{\ee^{sX}}Y=\ee^{sX}Ye^{-sX}.
\]
表示保持群乘法，因此
\begin{align*}
D_R(\ee^{sX})D_R(\ee^{tY})D_R(\ee^{-sX})
&=D_R\!\left(\ee^{sX}\ee^{tY}\ee^{-sX}\right).
\end{align*}
先对 $t$ 在 $t=0$ 求导，得到
\begin{align*}
D_R(\ee^{sX})\,dD_R(Y)\,D_R(\ee^{-sX})
&=dD_R\!\left(\operatorname{Ad}_{\ee^{sX}}Y\right).
\end{align*}
再对 $s$ 在 $s=0$ 求导。左边给矩阵对易子，右边给抽象伴随作用的导数：
\begin{align*}
[dD_R(X),dD_R(Y)]
&=dD_R([X,Y]_{\mathfrak g}).
\end{align*}
现在令 $X=X_a$、$Y=X_b$，并代入正文\eqnrefs{eq:lie-algebra-general-dp11}：
\begin{align*}
[dD_R(X_a),dD_R(X_b)]
&=f_{ab}{}^c\,dD_R(X_c).
\end{align*}
对酉表示，正文已定义 $dD_R(X_a)=-\ii T_R^a$，于是
\begin{align*}
[-\ii T_R^a,-\ii T_R^b]
&=f_{ab}{}^c(-\ii T_R^c),\\
-[T_R^a,T_R^b]
&=-\ii f_{ab}{}^cT_R^c.
\end{align*}
两边乘以 $-1$，得到正文\eqnrefs{eq:rep-lie-algebra-dp11}：
\[
[T_R^a,T_R^b]=\ii f_{ab}{}^cT_R^c.
\]
\end{derivation}

% DeepPass21: detailed Lie-algebra derivations gathered at the end of the owner subsection.










\begin{derivation}[从\eqnrefs{eq:generator-trace-normalization-dp11}到\eqnrefs{eq:suN-fierz-math-dp68}]



正文已经固定基本表示生成元的无迹性和迹归一化。四指标生成元和只能由两个 Kronecker 张量构成；用两种收缩固定它们的系数，就得到正文\eqnrefs{eq:suN-fierz-math-dp68}。设 $SU(N)$ 生成元满足
\begin{equation*}
 \Tr(T^aT^b)=T_F\delta^{ab},
 \qquad \Tr T^a=0.
\end{equation*}
把所有无迹厄米生成元与单位矩阵一起看成 $N\times N$ 矩阵空间的一组正交基。定义
\begin{equation*}
 X_{ij;kl}\equiv\sum_{a=1}^{N^2-1}(T^a)_{ij}(T^a)_{kl}.
\end{equation*}
由于 $i,l$ 与 $j,k$ 分别形成基本表示--反基本表示指标对，$SU(N)$ 不变量只能由
\begin{equation*}
 \delta_{il}\delta_{jk},
 \qquad
 \delta_{ij}\delta_{kl}
\end{equation*}
组成，所以写
\begin{equation*}
 X_{ij;kl}=A\delta_{il}\delta_{jk}+B\delta_{ij}\delta_{kl}.
\end{equation*}

第一组收缩取 $i=j$ 并求和：
\begin{align*}
 \sum_iX_{ii;kl}
 &=\sum_a\Tr(T^a)(T^a)_{kl}=0,\\
 &=A\sum_i\delta_{il}\delta_{ik}
 +B\sum_i\delta_{ii}\delta_{kl},\\
 &=(A+NB)\delta_{kl}.
\end{align*}
所以
\begin{equation*}
 B=-\frac AN.
\end{equation*}
第二组收缩取 $i=l$、$j=k$ 并对 $i,j$ 求和：
\begin{align*}
 \sum_{ij}X_{ij;ji}
 &=\sum_a\Tr(T^aT^a),\\
 &=T_F(N^2-1).
\end{align*}
而拟设右边给
\begin{align*}
 \sum_{ij}
 \left(A\delta_{ii}\delta_{jj}+B\delta_{ij}\delta_{ji}\right)
 &=AN^2+BN,\\
 &=A(N^2-1).
\end{align*}
故 $A=T_F$，从而
\begin{equation*}
 {
 \sum_a(T^a)_{ij}(T^a)_{kl}
 =T_F\left(
 \delta_{il}\delta_{jk}
 -\frac1N\delta_{ij}\delta_{kl}
 \right).}
\end{equation*}

检查：令 $k=j$ 并求和 $j$，左边是 $(\sum_aT^aT^a)_{il}=C_F\delta_{il}$；右边给
\begin{equation*}
 T_F\left(N-\frac1N\right)\delta_{il}
 =T_F\frac{N^2-1}{N}\delta_{il},
\end{equation*}
所以 $C_F=T_F(N^2-1)/N$，对 $T_F=1/2$ 恢复 $C_F=(N^2-1)/(2N)$。
\end{derivation}

\begin{derivation}[从\eqnrefs{eq:suN-fierz-math-dp68}到\eqnrefs{eq:ca-general-math-dp11}]



伴随 Casimir 可以不写任何显式 $SU(N)$ 矩阵，而由刚得到的完备关系求出。比较同一个对易子迹的两种计算即可固定 $C_A$。从基本对易子
\begin{equation*}
 [T^a,T^b]=\ii f^{abc}T^c
\end{equation*}
出发。考虑
\begin{equation*}
 \Tr([T^a,T^c][T^b,T^c]),
\end{equation*}
并对重复的 $c$ 求和。一方面直接代入结构常数：
\begin{align*}
 \sum_c\Tr([T^a,T^c][T^b,T^c])
 &=\sum_c\Tr(\ii f^{acd}T^d\,\ii f^{bce}T^e),\\
 &=-T_F f^{acd}f^{bcd},\\
 &=-T_F C_A\delta^{ab}.
\end{align*}
另一方面展开对易子：
\begin{align*}
 \sum_c\Tr([T^a,T^c][T^b,T^c])
 ={}&\sum_c\operatorname{Tr}\!\left[
 T^aT^cT^bT^c-T^aT^cT^cT^b\right.\\
 &\left.\hspace{2.5cm}-T^cT^aT^bT^c+T^cT^aT^cT^b
 \right].
\end{align*}
逐项用 Fierz/完备关系。对任意矩阵 $X$，由 Fierz 恒等式可得
\begin{equation*}
 \sum_c T^cXT^c
 =T_F\left[(\Tr X)I-\frac1N X\right].
\end{equation*}
因此
\begin{align*}
 \sum_c\Tr(T^aT^cT^bT^c)
 &=-\frac{T_F}{N}\Tr(T^aT^b),\\
 \sum_c\Tr(T^cT^aT^cT^b)
 &=-\frac{T_F}{N}\Tr(T^aT^b).
\end{align*}
另外 $\sum_cT^cT^c=C_FI$，故中间两项各为
\begin{equation*}
 -C_F\Tr(T^aT^b).
\end{equation*}
合并：
\begin{align*}
 \sum_c\Tr([T^a,T^c][T^b,T^c])
 &=-2\left(C_F+\frac{T_F}{N}\right)
 \Tr(T^aT^b),\\
 &=-2T_F\left(C_F+\frac{T_F}{N}\right)\delta^{ab}.
\end{align*}
取 $T_F=1/2$、$C_F=(N^2-1)/(2N)$：
\begin{equation*}
 2\left(C_F+\frac{T_F}{N}\right)
 =2\left(\frac{N^2-1}{2N}+\frac1{2N}\right)=N.
\end{equation*}
与第一种计算比较即得
\begin{equation*}
 C_A=N.
\end{equation*}
这条推导清楚展示了伴随 Casimir 与基本完备关系之间的关系。
\end{derivation}






















% DeepPass68: detailed differential-geometry derivations.

\begin{derivation}[从\eqnrefs{eq:exterior-derivative-definition-dp68}到\eqnrefs{eq:exterior-calculus-identities-dp68}]
外微分的两个基本性质都来自同一件事：普通偏导对光滑函数可以交换，而基一形式的楔积在交换次序时变号。对 $p$ 阶形式 $\omega$ 再作用一次外微分，
\begin{equation*}
 \dd^2 \omega
 =\frac1{p!}\partial_\rho\partial_\nu\omega_{\mu_1\cdots\mu_p}
 \dd x^\rho\wedge\dd x^\nu\wedge
 \dd x^{\mu_1}\wedge\cdots .
\end{equation*}
只保留 $(\rho,\nu)$ 的反对称部分即可，因为楔积会消去对称部分：
\begin{equation*}
 \partial_\rho\partial_\nu\omega\,
 \dd x^\rho\wedge\dd x^\nu
 =\frac12(\partial_\rho\partial_\nu-\partial_\nu\partial_\rho)\omega\,
 \dd x^\rho\wedge\dd x^\nu=0.
\end{equation*}
所以 $\dd^2=0$。

再令 $\alpha^{(p)}$ 为 $p$ 阶形式。对 $\alpha^{(p)}\wedge\beta$ 的系数求导时，导数落在 $\beta$ 上所产生的新基一形式必须穿过 $p$ 个属于 $\alpha$ 的基一形式，因此得到 $(-1)^p$：
\begin{equation*}
 \dd(\alpha^{(p)}\wedge\beta)
 =(\dd\alpha^{(p)})\wedge\beta
 +(-1)^p\alpha^{(p)}\wedge\dd\beta.
\end{equation*}
两条结果合在一起即为正文\eqnrefs{eq:exterior-calculus-identities-dp68}。
\end{derivation}

\begin{derivation}[从\eqnrefs{eq:connection-oneform-dp11}到\eqnrefs{eq:curvature-components-dp68}]
对表示空间取值的零形式 $\Psi$ 连续作用两次协变微分，
\begin{align*}
 D^2\Psi
 &=\dd(\dd\Psi+\ii g_{\rm G}\mathscr C\Psi)
 +\ii g_{\rm G}\mathscr C\wedge(\dd\Psi+\ii g_{\rm G}\mathscr C\Psi).
\end{align*}
用正文\eqnrefs{eq:exterior-calculus-identities-dp68} 中的 $\dd^2=0$ 和一形式的分次 Leibniz 规则，
\begin{equation*}
 \dd(\mathscr C\Psi)=(\dd\mathscr C)\Psi-\mathscr C\wedge\dd\Psi,
\end{equation*}
于是两个含 $\mathscr C\wedge\dd\Psi$ 的交叉项相消，得到
\begin{align*}
 D^2\Psi
 &=\ii g_{\rm G}(\dd\mathscr C)\Psi
 -g_{\rm G}^2\mathscr C\wedge\mathscr C\,\Psi\\
 &=\ii g_{\rm G}
 \left(\dd\mathscr C+\ii g_{\rm G}\mathscr C\wedge\mathscr C\right)\Psi.
\end{align*}
这就给出正文\eqnrefs{eq:curvature-form-dp11}。

再分别展开 $\dd\mathscr C$ 与 $\mathscr C\wedge\mathscr C$。首先
\begin{align*}
 \dd\mathscr C
 &=(\partial_\mu\mathscr C_\nu)\dd x^\mu\wedge\dd x^\nu\\
 &=\frac12(\partial_\mu\mathscr C_\nu-\partial_\nu\mathscr C_\mu)
 \dd x^\mu\wedge\dd x^\nu.
\end{align*}
其次，矩阵系数不能交换，
\begin{align*}
 \mathscr C\wedge\mathscr C
 &=\mathscr C_\mu\mathscr C_\nu\dd x^\mu\wedge\dd x^\nu\\
 &=\frac12[\mathscr C_\mu,\mathscr C_\nu]\dd x^\mu\wedge\dd x^\nu.
\end{align*}
与 $\mathcal F=\tfrac12\mathcal F_{\mu\nu}\dd x^\mu\wedge\dd x^\nu$ 比较系数，得到正文\eqnrefs{eq:curvature-components-dp68}。
\end{derivation}

\begin{derivation}[从\eqnrefs{eq:covariant-derivative-covariance-dp68}到\eqnrefs{eq:curvature-gauge-covariance-dp68}]
把 $\Psi'=U\Psi$ 代入正文\eqnrefs{eq:covariant-derivative-covariance-dp68}，左边为
\begin{align*}
 D'\Psi'
 &=\dd(U\Psi)+\ii g_{\rm G}\mathscr C'U\Psi\\
 &=(\dd U)\Psi+U\dd\Psi+\ii g_{\rm G}\mathscr C'U\Psi,
\end{align*}
而右边为
\begin{equation*}
 U(D\Psi)=U\dd\Psi+\ii g_{\rm G}U\mathscr C\Psi.
\end{equation*}
消去共同的 $U\dd\Psi$，并利用关系对任意 $\Psi$ 成立，
\begin{equation*}
 \dd U+\ii g_{\rm G}\mathscr C'U=\ii g_{\rm G}U\mathscr C.
\end{equation*}
右乘 $U^{-1}$ 后立即得到正文\eqnrefs{eq:finite-gauge-connection-dp11}。

由刚得到的联络变换式，协变导数满足
\[
D'=UDU^{-1}.
\]
为避免把算符乘法与普通矩阵乘法混淆，让两边作用在任意试探场 $\Psi'$ 上，并写 $\Psi'=U\Psi$。连续作用两次协变导数得到
\begin{align*}
D'^2\Psi'
&=D'\!\left(U D\Psi\right)\\
&=U D(D\Psi)\\
&=U D^2\Psi.
\end{align*}
正文对曲率的定义给出
\[
D^2\Psi=\ii g_{\rm G}\mathcal F\Psi,
\qquad
D'^2\Psi'=\ii g_{\rm G}\mathcal F'\Psi'.
\]
代入 $\Psi'=U\Psi$ 后，两种写法分别为
\begin{align*}
D'^2\Psi'
&=\ii g_{\rm G}\mathcal F' U\Psi,\\
D'^2\Psi'
&=\ii g_{\rm G}U\mathcal F\Psi.
\end{align*}
因为等式对任意 $\Psi$ 成立，且 $U$ 可逆，所以
\[
\mathcal F'U=U\mathcal F,
\qquad
\mathcal F'=U\mathcal F U^{-1},
\]
即正文\eqnrefs{eq:curvature-gauge-covariance-dp68}。
\end{derivation}

\begin{derivation}[从\eqnrefs{eq:covariant-commutator-curvature-dp68}到\eqnrefs{eq:nonabelian-bianchi-dp11}]
Bianchi 恒等式由协变导数对易子的 Jacobi 恒等式直接推出。写出
\begin{equation*}
 [D_\lambda,[D_\mu,D_\nu]]
 +[D_\mu,[D_\nu,D_\lambda]]
 +[D_\nu,[D_\lambda,D_\mu]]=0.
\end{equation*}
由正文\eqnrefs{eq:covariant-commutator-curvature-dp68}，第一项作用在任意 $\Psi$ 上为
\begin{align*}
 [D_\lambda,[D_\mu,D_\nu]]\Psi
 &=\ii g_{\rm G}[D_\lambda,\mathcal F_{\mu\nu}]\Psi\\
 &=\ii g_{\rm G}(D_\lambda\mathcal F_{\mu\nu})\Psi,
\end{align*}
其中
\begin{equation*}
 D_\lambda\mathcal F_{\mu\nu}
 =\partial_\lambda\mathcal F_{\mu\nu}
 +\ii g_{\rm G}[\mathscr C_\lambda,\mathcal F_{\mu\nu}].
\end{equation*}
对另外两项循环置换指标，Jacobi 恒等式便化为
\begin{equation*}
 \ii g_{\rm G}
 (D_\lambda\mathcal F_{\mu\nu}
 +D_\mu\mathcal F_{\nu\lambda}
 +D_\nu\mathcal F_{\lambda\mu})\Psi=0.
\end{equation*}
因为对任意 $\Psi$ 都成立，括号必须为零，即正文\eqnrefs{eq:nonabelian-bianchi-dp11}。
\end{derivation}

\begin{derivation}[从\eqnrefs{eq:parallel-transport-ode-dp68}到\eqnrefs{eq:wilson-endpoint-covariance-dp68}]
记沿路径的矩阵值系数为
\begin{equation*}
 A(s)=\dot x^\mu(s)\mathscr C_\mu[x(s)].
\end{equation*}
把正文\eqnrefs{eq:parallel-transport-ode-dp68} 从 $s_i$ 积到 $s$，得到 Volterra 型积分方程
\begin{equation*}
 \Psi(s)=\Psi(s_i)-\ii g_{\rm G}\int_{s_i}^{s}\dd s_1\,A(s_1)\Psi(s_1).
\end{equation*}
把右边的 $\Psi(s_1)$ 反复代回同一方程，前几阶为
\begin{align*}
 \Psi(s)=\Bigg[I
 &-\ii g_{\rm G}\int_{s_i}^{s}\dd s_1A(s_1)\\
 &+(-\ii g_{\rm G})^2
 \int_{s_i}^{s}\dd s_1\int_{s_i}^{s_1}\dd s_2\,
 A(s_1)A(s_2)\\
 &+(-\ii g_{\rm G})^3
 \int_{s_i}^{s}\dd s_1\int_{s_i}^{s_1}\dd s_2
 \int_{s_i}^{s_2}\dd s_3\,
 A(s_1)A(s_2)A(s_3)+\cdots\Bigg]\Psi(s_i).
\end{align*}
第 $n$ 阶积分域满足 $s_1\ge s_2\ge\cdots\ge s_n$；由于不同 $s$ 处的矩阵一般不交换，较大的路径参数必须排在乘积左侧。用 $\mathcal P$ 记这种顺序，级数重求和为
\begin{equation*}
 \Psi(x_f)=U_C\Psi(x_i),\qquad
 U_C=\mathcal P\exp\!\left[-\ii g_{\rm G}\int_C\mathscr C_\mu\dd x^\mu\right],
\end{equation*}
即正文\eqnrefs{eq:wilson-line-general-dp11}。

现在改变内部空间的局域基底，令 $\Psi'(x)=U(x)\Psi(x)$。同一个输运关系在新基底中满足
\begin{align*}
 \Psi'(x_f)
 &=U(x_f)\Psi(x_f)\\
 &=U(x_f)U_C\Psi(x_i)\\
 &=U(x_f)U_CU^{-1}(x_i)\Psi'(x_i).
\end{align*}
另一方面，新基底中的 Wilson 线 $U_C'$ 按正文定义满足
\begin{equation*}
 \Psi'(x_f)=U_C'\Psi'(x_i).
\end{equation*}
因为该等式对任意初始内部向量 $\Psi'(x_i)$ 都成立，比较线性算符的系数得到
\begin{equation*}
 U_C'=U(x_f)U_CU^{-1}(x_i),
\end{equation*}
即正文\eqnrefs{eq:wilson-endpoint-covariance-dp68}。
\end{derivation}

\begin{derivation}[从\eqnrefs{eq:wilson-line-general-dp11}到\eqnrefs{eq:small-loop-curvature-dp11}]
取左下角为 $x$、两条边向量分别为 $a^\mu$ 与 $b^\nu$ 的小矩形，只保留面积阶 $ab$。四条边的短程 Wilson 线依次为
\begin{align*}
 U_1&=I-\ii g_{\rm G}\mathscr C_\mu(x)a^\mu+\order(a^2),\\
 U_2&=I-\ii g_{\rm G}[\mathscr C_\nu(x)+a^\mu\partial_\mu\mathscr C_\nu(x)]b^\nu
 +\order(b^2,a^2b),\\
 U_3&=I+\ii g_{\rm G}[\mathscr C_\mu(x)+b^\nu\partial_\nu\mathscr C_\mu(x)]a^\mu
 +\order(a^2,ab^2),\\
 U_4&=I+\ii g_{\rm G}\mathscr C_\nu(x)b^\nu+\order(b^2).
\end{align*}
闭合输运为 $U_\square=U_4U_3U_2U_1$。所有一阶项相消；Taylor 项留下
\begin{equation*}
 -\ii g_{\rm G}(\partial_\mu\mathscr C_\nu-\partial_\nu\mathscr C_\mu)a^\mu b^\nu,
\end{equation*}
矩阵乘法的二阶交叉项留下
\begin{equation*}
 -g_{\rm G}^2[\mathscr C_\mu,\mathscr C_\nu]a^\mu b^\nu.
\end{equation*}
合并并使用正文\eqnrefs{eq:curvature-components-dp68}，便得到正文\eqnrefs{eq:small-loop-curvature-dp11}。
\end{derivation}

\begin{derivation}[从\eqnrefs{eq:pontryagin-form-definition-dp68}到\eqnrefs{eq:pontryagin-number-dp11}]
以下计算先核对两种写法的归一化；Pontryagin 数的整数性来自整体拓扑条件，而不是这一局部代数恒等式。写
\begin{equation*}
 \mathcal F=\frac12\mathcal F_{\mu\nu}^aT^a
 \dd x^\mu\wedge\dd x^\nu,
 \qquad
 \Tr(T^aT^b)=\frac12\delta^{ab}.
\end{equation*}
于是
\begin{align*}
 \Tr(\mathcal F\wedge\mathcal F)
 &=\frac18\mathcal F_{\mu\nu}^a\mathcal F_{\rho\sigma}^a
 \epsilon^{\mu\nu\rho\sigma}\dd^4 x.
\end{align*}
再用对偶场强定义
\begin{equation*}
 {}^\star\!\mathcal F^{a\mu\nu}
 =\frac12\epsilon^{\mu\nu\rho\sigma}\mathcal F_{\rho\sigma}^a,
\end{equation*}
可写成
\begin{equation*}
 \Tr(\mathcal F\wedge\mathcal F)
 =\frac14\mathcal F_{\mu\nu}^a
 {}^\star\!\mathcal F^{a\mu\nu}\dd^4 x.
\end{equation*}
代回正文\eqnrefs{eq:pontryagin-form-definition-dp68}，得到\eqnrefs{eq:pontryagin-number-dp11} 中的分量归一化 $g_{\rm G}^2/(32\pi^2)$。其中 $Q_{\rm top}\in\mathbb Z$ 的部分依赖正文已经说明的边界条件与整体拓扑定理，不属于这个代数推导的结论。
\end{derivation}


\section{散射积分与长时间极限}

连续积分核、带源二次积分和函数积分的有限维代数确定以后，散射计算还需要把多个分母、能量守恒和有限作用时间转化为可积表达式。参数积分与长时间分布极限正是从形式积分走向有限截面与跃迁率的一组通用数学工具。

\subsection{参数积分与多维动量积分}
多维动量积分反复需要三种操作：用参数积分合并多个二次分母，把平移后的动量积分化成标准高斯型，以及在含 $\delta$ 函数的变量替换中保留正确的 Jacobian。标准高斯积分、Gamma 函数恒等式和 $\delta[f]$ 根公式可直接调用；真正需要逐步追踪的是参数化怎样改变分母、$\ii0^+$ 怎样固定解析边界，以及调节器是否保持平移与 Lorentz 对称性。有限作用时间造成的能量选择则由时间 Fourier 核单独控制。

% -----------------------------------------------------------------------------
% Common mathematical tools for generic QFT, proved before first use
% -----------------------------------------------------------------------------

散射积分反复遇到三个具体困难：能量守恒的 $\delta$ 函数需要在变量变换后给出正确 Jacobian；多个传播子分母需要合并成一个可平移动量的分母；圈动量的张量分子需要利用对称性化成标量积分。这三类工具共同固定散射积分中的 Jacobian、分母组合以及张量约化，并同时给出收敛域与解析延拓条件。

若光滑函数 $f(x)$ 的零点 $x_i$ 都是单根，即 $f(x_i)=0$ 且 $f'(x_i)\neq0$，则复合 $\delta$ 分布满足
\begin{equation}
\boxed{
\delta[f(x)]
=\sum_i\frac{\delta(x-x_i)}{|f'(x_i)|}.}
\label{eq:delta-composition-r68}
\end{equation}
这个 Jacobian 因子由变量变换唯一确定，并直接进入质量壳积分与末态相空间测度。

圈积分最常见的问题是多个传播子分母相乘。母公式从每个分母的 Schwinger 参数表示开始。若 $\operatorname{Re}A_i>0$、$\operatorname{Re}\nu_i>0$，则
\begin{equation}
 \boxed{
 \frac1{A_i^{\nu_i}}
 =\frac1{\Gamma(\nu_i)}
 \int_0^\infty\dd\tau_i\,\tau_i^{\nu_i-1}\ee^{-A_i\tau_i}.}
 \label{eq:schwinger-parameter-r9}
\end{equation}
令 $\nu\equiv\sum_{i=1}^N\nu_i$，并作变量分解 $\tau_i=sx_i$、$x_i\ge0$、$\sum_i x_i=1$。对公共尺度 $s$ 的积分给出 $\Gamma(\nu)$，于是任意 $N$ 个分母满足
\begin{equation}
\boxed{
 \frac1{\prod_{i=1}^N A_i^{\nu_i}}
 =\frac{\Gamma(\nu)}{\prod_{i=1}^N\Gamma(\nu_i)}
 \int_{x_i\ge0}\!\left[\prod_{i=1}^N\dd x_i\,x_i^{\nu_i-1}\right]
 \delta\!\left(1-\sum_{i=1}^N x_i\right)
 \frac1{\left(\sum_{i=1}^N x_iA_i\right)^{\nu}}.}
\label{eq:feynman-n-general-dp74}
\end{equation}
这就是后续圈图使用的母参数化；参数 $x_i$ 无量纲，Minkowski 传播子的 $\ii0^+$ 处方随线性组合一起保留。对 $N=2$，\eqnrefs{eq:feynman-n-general-dp74} 化为
\begin{equation}
\boxed{
\frac1{A^{\nu_A}B^{\nu_B}}
=\frac{\Gamma(\nu_A+\nu_B)}{\Gamma(\nu_A)\Gamma(\nu_B)}
\int_0^1\dd x\,
\frac{x^{\nu_A-1}(1-x)^{\nu_B-1}}
{[xA+(1-x)B]^{\nu_A+\nu_B}}.}
\label{eq:feynman-two-general-r68}
\end{equation}
再取 $\nu_A=\nu_B=1$ 才得到常用特例
\begin{equation}
 \boxed{
 \frac1{AB}
 =\int_0^1\dd x\,
 \frac1{[xA+(1-x)B]^2}.}
 \label{eq:feynman-parameter-proof-r9}
\end{equation}
对三个一次分母同样只是 $N=3$、$\nu_i=1$ 的特例：
\begin{equation}
 \boxed{
 \frac1{ABC}
 =2\int_0^1\dd x\dd y\dd z\,
 \delta(x+y+z-1)
 \frac1{(xA+yB+zC)^3}.}
 \label{eq:feynman-three-denominator-r14}
\end{equation}
前因子 $2=\Gamma(3)$ 来自母公式中的公共尺度积分，而不是某张三角圈图的组合因子。

完成多个二次分母的合并和平移积分变量后，许多 $d$ 维 Euclidean 动量积分都会归结为高斯矩。对 $a>0$，$d$ 维 Euclidean 向量 $\ell_E$ 满足
\begin{equation}
 \boxed{
 \int\dd^d \ell_E\,\ee^{-a\ell_E^2}
 =\left(\frac\pi a\right)^{d/2}.}
 \label{eq:d-dimensional-gaussian-r9}
\end{equation}
这里的 $d$ 暂时不要求等于物理时空的整数维数。某些量子场论动量积分在 $d=4$ 时发散，却在一段 $\operatorname{Re}d$ 范围内收敛。\emph{维数正规化}（dimensional regularization）把收敛维数范围内得到的解析表达式作为关于 $d$ 的函数，并将它解析延拓到复维数；物理四维由
\begin{equation*}
 d=4-2\epsilon,
 \qquad \epsilon\to0.
\end{equation*}
这样，当积分变量的模 $|\ell|\to\infty$ 时产生的高动量发散会表现为 $1/\epsilon$ 等极点；在量子场论中，这类来自任意大内部动量的发散称为\emph{紫外发散}（ultraviolet divergence）。有限部分仍由同一个解析函数给出。这里的“非整数维”是给发散积分定义解析延拓的调节参数，并不是说物理时空真的具有分数维数。维数正规化之所以特别适合相对论场论，是因为它可以在计算过程中保持 Lorentz 协变的张量结构；真正的四维物理量要在完成参数重定义和取 $\epsilon\to0$ 极限后得到。

若调节方案保持 $d$ 维 Lorentz 对称性，对只依赖 $\ell^2$ 的标量函数 $f$，二阶张量积分只能与度规成正比：
\begin{equation}
\int\dd^d \ell\,\ell^\mu\ell^\nu f(\ell^2)
=C_f\eta^{\mu\nu}.
\label{eq:rank2-isotropic-ansatz-dp68}
\end{equation}
与 $\eta_{\mu\nu}$ 收缩固定 $C_f$ 后得到
\begin{equation}
 \boxed{
 \int\dd^d \ell\,\ell^\mu\ell^\nu f(\ell^2)
 =\frac{\eta^{\mu\nu}}d
 \int\dd^d \ell\,\ell^2 f(\ell^2).}
 \label{eq:tensor-reduction-proof-r9}
\end{equation}
这一步依赖调节器保持平移和 Lorentz 对称性；对尚未调节的发散积分不能先做形式上的“对称平均”。

% DeepPass68: general d-dimensional integral identities for loop calculations.

$d$ 维径向积分需要单位球面的面积。由定义
\begin{equation}
\Omega_d\equiv\int\dd\Omega_{d-1},
\label{eq:d-sphere-solid-angle-definition-dp68}
\end{equation}
可得
\begin{equation}
\boxed{
\Omega_d=\frac{2\pi^{d/2}}{\Gamma(d/2)}.}
\label{eq:d-sphere-solid-angle-dp21}
\end{equation}

把 Schwinger 参数式用于 Euclidean 分母，先写成
\begin{equation}
\frac1{(\ell_E^2+\Delta)^\nu}
=\frac1{\Gamma(\nu)}
\int_0^\infty\dd s\,s^{\nu-1}\ee^{-s(\ell_E^2+\Delta)}.
\label{eq:euclidean-master-schwinger-dp68}
\end{equation}
在 $\Re\Delta>0$、$\Re\nu>d/2$ 的收敛域内先完成 $\ell_E$ 与 $s$ 积分，得到
\begin{equation}
\boxed{
\int\frac{\dd^d \ell_E}{(2\pi)^d}
\frac1{(\ell_E^2+\Delta)^\nu}
=\frac1{(4\pi)^{d/2}}
\frac{\Gamma(\nu-d/2)}{\Gamma(\nu)}
\Delta^{d/2-\nu}.}
\label{eq:euclidean-master-integral-dp21}
\end{equation}
右端随后定义对 $d$、$\nu$ 与 $\Delta$ 的解析延拓。

四阶张量分子需要比\eqnrefs{eq:tensor-reduction-proof-r9} 多一个对称结构。完全对称且各向同性的秩四张量只能写成
\begin{equation}
\int\dd^d \ell\,\ell^\mu\ell^\nu\ell^\rho\ell^\sigma f(\ell^2)
=C_4\bigl(
\eta^{\mu\nu}\eta^{\rho\sigma}
+\eta^{\mu\rho}\eta^{\nu\sigma}
+\eta^{\mu\sigma}\eta^{\nu\rho}\bigr).
\label{eq:rank4-isotropic-ansatz-dp68}
\end{equation}
双重收缩固定系数后得到
\begin{equation}
\boxed{
\int\dd^d \ell\,
\ell^\mu\ell^\nu\ell^\rho\ell^\sigma f(\ell^2)
=\frac{\eta^{\mu\nu}\eta^{\rho\sigma}
+\eta^{\mu\rho}\eta^{\nu\sigma}
+\eta^{\mu\sigma}\eta^{\nu\rho}}
{d(d+2)}
\int\dd^d \ell\,(\ell^2)^2f(\ell^2).}
\label{eq:rank4-tensor-reduction-dp21}
\end{equation}

在 $d=4-2\epsilon$ 附近，圈积分中经常出现同一个无量纲组合
\begin{equation}
\mathcal U_\epsilon(\Delta)
\equiv\Gamma(\epsilon)(4\pi)^\epsilon
\left(\frac{\Delta}{\mu^2}\right)^{-\epsilon},
\label{eq:dimreg-universal-factor-dp68}
\end{equation}
其中 $\mu$ 是保持重整化参数具有固定工程量纲所引入的尺度。小 $\epsilon$ 展开为
\begin{equation}
\boxed{
\mathcal U_\epsilon(\Delta)
=\frac1\epsilon-\gamma_E+\ln4\pi
-\ln\frac{\Delta}{\mu^2}+\order(\epsilon).}
\label{eq:dimreg-universal-expansion-dp68}
\end{equation}
因此 $\overline{\rm MS}$ 中反复出现的 $1/\epsilon-\gamma_E+\ln4\pi$ 只是这个通用展开的方案常数部分。


\subsection{长时间分布极限}
有限作用时间首先给出宽度约为 $1/T$ 的谱选择函数，而不是严格的能量 $\delta$ 函数。只有在对平滑末态谱密度积分并取 $T\to\infty$ 后，$\mathrm{sinc}^2$ 核才以分布的意义趋于能量守恒条件；这一分布极限正是散射率和 Fermi 黄金律中能量守恒的数学来源。

% Long-time distribution targets; detailed proof follows at the end of this structure.

有限观测时间不会产生严格的能量守恒 $\delta$ 函数，而产生宽度约为 $\hbar/T$ 的有限时间谱核。令 $\delta E$ 为实能量失配，单位 J；令 $T>0$ 为相互作用或观测时间，单位 s。定义
\begin{equation*}
 F_T(\delta E)
 \equiv T\,\operatorname{sinc}^2\!\left(\frac{\delta E\,T}{2\hbar}\right),
 \qquad
 \operatorname{sinc}x\equiv\frac{\sin x}{x}.
\end{equation*}
因为 $F_T$ 的峰高随 $T$ 增长而峰宽随 $1/T$ 缩小，它不能按普通逐点函数极限理解；正确极限必须用分布作用在测试函数上的方式定义。

对任意光滑、足够快衰减的测试函数 $g(\delta E)$，定义
\begin{equation*}
 I_T[g]\equiv
 \int_{-\infty}^{\infty}\dd(\delta E)\,
 F_T(\delta E)g(\delta E).
\end{equation*}
作无量纲变量替换
$x=\delta E\,T/(2\hbar)$，于是
\begin{equation}
 \boxed{
 I_T[g]
 =2\hbar\int_{-\infty}^{\infty}\dd x\,
 \left(\frac{\sin x}{x}\right)^2
 g\!\left(\frac{2\hbar x}{T}\right).}
 \label{eq:sinc-test-substitution-r15}
\end{equation}
因此长时间极限只剩两个问题：$g(2\hbar x/T)\to g(0)$，以及核 $\operatorname{sinc}^2x$ 的总面积是多少。

为不把 Dirichlet 积分作为未证公式，引入指数调节器 $\epsilon>0$，定义
\begin{equation*}
 D_\epsilon(b)
 \equiv\int_0^\infty\dd x\,\ee^{-\epsilon x}\frac{\sin bx}{x}.
\end{equation*}
对 $b$ 求导后积分变成初等 Laplace 变换，最终得到
\begin{equation}
 \boxed{
 D_\epsilon(b)=\arctan\frac{b}{\epsilon}
 \xrightarrow[\epsilon\to0^+]{}
 \frac{\pi}{2}\operatorname{sgn}b.}
 \label{eq:dirichlet-integral-derived-r15}
\end{equation}
由此进一步得到
$\int_{-\infty}^{\infty}(\sin x/x)^2\dd x=\pi$。代回\eqnrefs{eq:sinc-test-substitution-r15} 后，有限时间谱核的分布极限为
\begin{equation}
 \boxed{
 \lim_{T\to\infty}
 T\,\operatorname{sinc}^2\!\left(\frac{\delta E\,T}{2\hbar}\right)
 =2\pi\hbar\,\delta(\delta E).}
 \label{eq:sinc-delta-proved-r15}
\end{equation}
\eqnrefs{eq:sinc-delta-proved-r15} 的等号仅表示对测试函数积分后的分布收敛；任何有限 $T$ 的实际过程仍必须保留宽度有限的 $\operatorname{sinc}^2$ 线形，不能把它提前替换成严格 $\delta$ 函数。


傅里叶变换、$\delta/\Theta$ 分布、$\ii0^+$ 处方、连续态测度、算符伴随、投影/预解式、高斯/Grassmann 积分以及有限时间谱核共同固定了全书的数学底座。群论、电动力学与场量子化只在这些对象上增加新的物理结构，因而不同章节中的前因子、测度和边界处方可以直接比较。


% Detailed derivations gathered at the end of the owning structure.


